\documentclass[aip,jcp,reprint,noshowkeys,superscriptaddress]{revtex4-2}
\usepackage{graphicx,dcolumn,bm,xcolor,microtype,multirow,amscd,amsmath,amssymb,amsfonts,physics,wrapfig,bbold,siunitx,xspace,braket,txfonts}
\usepackage[version=4]{mhchem}
\let\openbox\relax % revtex defines \openbox; amsthm redefines it
\usepackage{amsthm}
\theoremstyle{plain}
\newtheorem{theorem}{Theorem}
\newtheorem{lemma}{Lemma}
\newtheorem{corollary}{Corollary}
\theoremstyle{definition}
\newtheorem{definition}{Definition}
\newtheorem{assumption}{Assumption}
\theoremstyle{remark}
\newtheorem{remark}{Remark}

% Folder of the main text, seen from where TeX runs: ../ when si.tex is built in anc/,
% empty on Overleaf, which runs TeX in the project root next to manuscript.tex.
\IfFileExists{manuscript.tex}{\def\maintextdir{}}{%
  \IfFileExists{manuscript/manuscript.tex}{\def\maintextdir{manuscript/}}{%
    \def\maintextdir{../}}}
\usepackage{xr-hyper}
\externaldocument{\maintextdir manuscript}

\usepackage[utf8]{inputenc}
\usepackage[T1]{fontenc}
\usepackage[english]{babel}
\babelprovide[import=en]{en} % bib entries carry language = {en}
\usepackage{feynmp-auto}
\usepackage{longtable}
\usepackage{rotating}
\usepackage{hyperref}
\usepackage[compat=1.0.0]{tikz-feynman}

\hypersetup{
    colorlinks,
    linkcolor={red!50!black},
    citecolor={red!70!black},
    urlcolor={red!80!black}
}

\usepackage{listings}
\definecolor{codegreen}{rgb}{0.58,0.4,0.2}
\definecolor{codegray}{rgb}{0.5,0.5,0.5}
\definecolor{codepurple}{rgb}{0.25,0.35,0.55}
\definecolor{codeblue}{rgb}{0.30,0.60,0.8}
\definecolor{backcolour}{rgb}{0.98,0.98,0.98}
\definecolor{mygray}{rgb}{0.5,0.5,0.5}

\definecolor{sqred}{rgb}{0.85,0.1,0.1}
\definecolor{sqgreen}{rgb}{0.25,0.65,0.15}
\definecolor{sqorange}{rgb}{0.90,0.50,0.15}
\definecolor{sqblue}{rgb}{0.10,0.3,0.60}

\lstdefinestyle{mystyle}{
    backgroundcolor=\color{backcolour},
    commentstyle=\color{codegreen},
    keywordstyle=\color{codeblue},
    numberstyle=\tiny\color{codegray},
    stringstyle=\color{codepurple},
    basicstyle=\ttfamily\footnotesize,
    breakatwhitespace=false,
    breaklines=true,
    captionpos=b,
    keepspaces=true,
    numbers=left,
    numbersep=5pt,
    numberstyle=\ttfamily\tiny\color{mygray},
    showspaces=false,
    showstringspaces=false,
    showtabs=false,
    tabsize=2
  }

  \newcolumntype{d}{D{.}{.}{-1}}

  \lstset{style=mystyle}

\newcommand{\cre}[1]{a_{#1}^\dagger}
\newcommand{\ani}[1]{a_{#1}} 

\newcommand{\fcre}[1]{\Hat{a}_{#1}^{\dag}}
\newcommand{\fani}[1]{\Hat{a}_{#1}}

\newcommand{\bfcre}[1]{\Hat{\Bar{a}}_{#1}^{\dag}}
\newcommand{\bfani}[1]{\Hat{\Bar{a}}_{#1}}

\newcommand{\bcre}[1]{\Hat{b}_{#1}^{\dag}}
\newcommand{\bani}[1]{\Hat{b}_{#1}}

\newcommand{\bacre}[1]{\Bar{b}_{#1}^{\dag}}
\newcommand{\baani}[1]{\Bar{b}_{#1}}

\newcommand{\bbacre}[1]{\dBar{b}_{#1}^{\dag}}
\newcommand{\bbaani}[1]{\dBar{b}_{#1}}

\newcommand{\bbcre}[1]{\Hat{\beta}_{#1}^{\dag}}
\newcommand{\bbani}[1]{\Hat{\beta}_{#1}}


\newcommand{\ccre}[1]{\Hat{c}_{#1}^{\dag}}
\newcommand{\cani}[1]{\Hat{c}_{#1}}

\newcommand{\bccre}[1]{\Hat{\Bar{c}}_{#1}^{\dag}}
\newcommand{\bcani}[1]{\Hat{\Bar{c}}_{#1}}


\newcommand{\dcre}[1]{\Hat{d}_{#1}^{\dag}}
\newcommand{\dani}[1]{\Hat{d}_{#1}}

\newcommand{\bdcre}[1]{\Hat{\Bar{d}}_{#1}^{\dag}}
\newcommand{\bdani}[1]{\Hat{\Bar{d}}_{#1}}

\usepackage{accents}
\newcommand{\dBar}[1]{\Bar{\Bar{#1}}}

\newcommand{\NO}[1]{\{#1\}} 
\newcommand{\ERI}[2]{\braket{#1|#2}}
\newcommand{\sERI}[2]{(#1|#2)}
\newcommand{\dbERI}[2]{\mel{#1}{}{#2}}
\newcommand{\e}[2]{\epsilon_{#1}^{#2}}
\newcommand{\eHF}[1]{\epsilon_{#1}^\text{HF}}
\newcommand{\eKS}[1]{\epsilon_{#1}^\text{KS}}
\newcommand{\Om}{\Omega}
\newcommand{\om}{\omega}

\newcommand{\titou}[1]{\textcolor{red}{#1}}
\newcommand{\petros}[1]{\textcolor{cyan}{#1}}
\newcommand{\jt}[1]{\textcolor{orange}{#1}}
\newcommand{\ag}[1]{\textcolor{teal}{#1}}
\newcommand{\ai}[1]{\textcolor{magenta}{#1}}
\newcommand{\trashPFL}[1]{\textcolor{red}{\sout{#1}}}
\newcommand{\trashMPK}[1]{\textcolor{cyan}{\sout{#1}}}
\newcommand{\PFL}[1]{\titou{(\underline{\bf PFL}: #1)}}
\newcommand{\MPK}[1]{\petros{(\underline{\bf MPK}: #1)}}
\newcommand{\JT}[1]{\jt{(\underline{\bf JT}: #1)}}
\newcommand{\AG}[1]{\ag{(\underline{\bf AG}: #1)}}
\newcommand{\AI}[1]{\ai{(\underline{\bf AI}: #1)}}
\newcommand{\eqrref}[1]{Eq.~\eqref{#1}}
\newcommand{\fnm}{\footnotemark}
\newcommand{\fnt}{\footnotetext}

\newcommand{\HF}{\text{HF}}
\newcommand{\GW}{\text{$GW$}}
\newcommand{\GOWO}{\text{$G_0W_0$}}
\newcommand{\drCCD}{\text{drCCD}}
\newcommand{\RPA}{\text{RPA}}
\newcommand{\IP}{\text{IP}}
\newcommand{\EA}{\text{EA}}
\newcommand{\IPEA}{\text{IP/EA}}
\newcommand{\EOM}{\text{EOM}}

\newcommand{\TCC}{\text{TCC}}
\newcommand{\ECC}{\text{ECC}}

\newcommand{\ba}{\boldsymbol{a}}
\newcommand{\bb}{\boldsymbol{b}}
\newcommand{\bbb}{\boldsymbol{\beta}}
\newcommand{\bc}{\boldsymbol{c}}
\newcommand{\bbc}{\Bar{\bc}}
\newcommand{\bd}{\boldsymbol{d}}
\newcommand{\bbd}{\Bar{\bd}}
\newcommand{\bx}{\boldsymbol{x}}
\newcommand{\bA}{\boldsymbol{A}}
\newcommand{\bB}{\boldsymbol{B}}
\newcommand{\bC}{\boldsymbol{C}}
\newcommand{\bD}{\boldsymbol{D}}
\newcommand{\bG}{\boldsymbol{G}}
\newcommand{\bX}{\boldsymbol{X}}
\newcommand{\bY}{\boldsymbol{Y}}
\newcommand{\bT}{\boldsymbol{T}}
\newcommand{\bt}{\boldsymbol{t}}
\newcommand{\bR}{\boldsymbol{R}}
\newcommand{\br}{\boldsymbol{r}}
\newcommand{\bE}{\boldsymbol{E}}
\newcommand{\bM}{\boldsymbol{M}}
\newcommand{\bN}{\boldsymbol{N}}
\newcommand{\bH}{\boldsymbol{H}}
\newcommand{\bO}{\boldsymbol{0}}
\newcommand{\bI}{\boldsymbol{1}}
\newcommand{\bS}{\boldsymbol{S}}
\newcommand{\bz}{\boldsymbol{z}}
\newcommand{\bV}{\boldsymbol{V}}
\newcommand{\bff}{\boldsymbol{f}}
\newcommand{\bLam}{\boldsymbol{\Lambda}}
\newcommand{\blam}{\boldsymbol{\lambda}}
\newcommand{\bgam}{\boldsymbol{\gamma}}
\newcommand{\bOm}{\boldsymbol{\Omega}}
\newcommand{\bSig}{\boldsymbol{\Sigma}}
\newcommand{\beps}{\boldsymbol{\epsilon}}

\newcommand{\hT}{\Hat{T}}
\newcommand{\hW}{\Hat{W}}
\newcommand{\hF}{\Hat{F}}
\newcommand{\hL}{\Hat{L}}
\newcommand{\hR}{\Hat{R}}
\newcommand{\hH}{\Hat{H}}
\newcommand{\hbH}{\Bar{H}}
\newcommand{\hbbH}{\dBar{H}}
\newcommand{\hI}{\Hat{1}}
\newcommand{\hV}{\Hat{V}}
\newcommand{\hZ}{\Hat{Z}}
\newcommand{\hLam}{\Hat{\Lambda}}
\newcommand{\htau}{\Hat{\tau}}
\newcommand{\hXi}{\Hat{\Xi}}

\newcommand{\cL}{\mathcal{L}}
\newcommand{\wtH}{\widetilde{H}}
\newcommand{\wtV}{\widetilde{V}}

\newcommand{\T}[1]{#1^{\intercal}}
\newcommand{\Ec}{E_\text{c}}

\newcommand{\ldr}{\text{$\lambda$-dr}}
\newcommand{\dr}{\text{dr}}
\newcommand{\ii}{\text{i}}
\newcommand{\SupInf}{\textcolor{blue}{Supporting Information}\xspace}

\newcommand{\stkout}[1]{\ifmmode\text{\sout{\ensuremath{#1}}}\else\sout{#1}\fi}

\renewcommand{\tableautorefname}{Table}
\renewcommand{\sectionautorefname}{Sec.}
\renewcommand{\figureautorefname}{Fig.}
\renewcommand\thefigure{S\arabic{figure}}
\renewcommand\thetable{S\arabic{table}}

% addresses
\newcommand{\UnHam}{Department of Chemistry, University of Hamburg, 22761 Hamburg, Germany; The Hamburg Centre for Ultrafast Imaging (CUI), Hamburg 22761, Germany}
\newcommand{\LCPQ}{Laboratoire de Chimie et Physique Quantiques (UMR 5626), Universit\'e de Toulouse, CNRS, Toulouse, France}
\newcommand{\VUA}{Department of Chemistry and Pharmaceutical Sciences, Vrije Universiteit, De Boelelaan 1105, 1081 HV Amsterdam, The Netherlands}


\begin{document}	

\title{Perspective: Consistent resummation schemes for the one-particle Green's function}

\author{Johannes T\"olle$^{\ast}$}
        \email{johannes.toelle@uni-hamburg.de}
        \affiliation{\UnHam}

\author{Arno Förster}
	\email{a.t.l.foerster@vu.nl}
	\affiliation{\VUA}

\author{Pierre-Fran\c{c}ois \surname{Loos}$^{*}$}
	\email{loos@irsamc.ups-tlse.fr}
	\affiliation{\LCPQ}


\maketitle

\tableofcontents

\section{\label{sec:si:adc}ADC constructions}
\subsection{Notation}
We work with the canonical Hartree--Fock (HF) spin orbitals of an $N$-electron reference determinant $\ket{\Phi_0}$, with orbital energies $\epsilon_p$ and real orbitals. Indices $i,j,k,l$ denote occupied, $a,b,c,d$ virtual and $p,q,r,s$ arbitrary spin orbitals. $\ket{\Psi_0}$ and $E_0$ are the exact ground state and its energy, and we assume that every virtual orbital energy lies above every occupied one.

\subsection{General setting and the union space}
The goal of the Dyson algebraic diagrammatic construction is to construct an effective Hamiltonian that combines the operators $H^{N+1}-E_0$ and $H^{N-1}-E_0$, which individually describe electron attachment and removal,
\begin{equation}
\begin{aligned}
\left(H^{N+1}-E_0\right)\ket{\Psi^{N+1}_n} ={}& (E^{N+1}_n-E_0)\ket{\Psi^{N+1}_n} \;,\\
-\left(H^{N-1}-E_0\right)\ket{\Psi^{N-1}_k} ={}& (E_0-E^{N-1}_k)\ket{\Psi^{N-1}_k} \;.
\end{aligned}
\end{equation}
The direct sum of these operators,
\begin{equation}
\mathcal{H}=\big(H^{N+1}-E_0\big)\oplus\big(-(H^{N-1}-E_0)\big) \;,
\end{equation}
acts on the direct sum $\mathcal{V}$ of the $(N+1)$- and $(N-1)$-electron Hilbert spaces as
\begin{equation}
\begin{aligned}
\mathcal{H}\,\big(\ket{\Psi^{N+1}_n};0\big)={}&(E^{N+1}_n-E_0)\big(\ket{\Psi^{N+1}_n};0\big) \;,\\
\mathcal{H}\,\big(0;\ket{\Psi^{N-1}_k}\big)={}&(E_0-E^{N-1}_k)\big(0;\ket{\Psi^{N-1}_k}\big) \;.
\end{aligned}
\end{equation}
We refer to $\mathcal{V}$ as the \emph{union space}. It is spanned by vectors generated from the exact ground state by operators $Y$ that add one electron,
\begin{equation}
v(Y) = \big(Y\ket{\Psi_0}\,;\,Y^\dagger\ket{\Psi_0}\big) \;.
\label{eq:si:vY}
\end{equation}
The upper component lies in the $(N+1)$-electron space. Because of the Hermitian conjugation, the lower component lies in the $(N-1)$-electron space. We order these vectors in classes. Class 0 holds the orbital vectors
\begin{equation}
v(a_p^\dagger) = \big(a_p^\dagger\ket{\Psi_0}\,;\,a_p\ket{\Psi_0}\big)
\label{eq:si:class0}
\end{equation}
for all orbitals $p$. Class 1 holds the vectors generated by the 2p1h strings
\begin{equation}
O^{+}_\nu = a_a^\dagger a_b^\dagger a_k \;,
\label{eq:si:class1plus}
\end{equation}
and by the adjoints of the 2h1p strings,
\begin{equation}
O^{-\dagger}_\mu = a_i^\dagger a_j^\dagger a_a \;,
\label{eq:si:class1minus}
\end{equation}
whose lower component is the 2h1p state
\begin{equation}
O^{-}_\mu\ket{\Psi_0} = a_a^\dagger a_j a_i\ket{\Psi_0} \;.
\end{equation}
Class $k$ holds, in the same way, the vectors generated by the $(k+1)$p$k$h strings and by the adjoints of the $(k+1)$h$k$p strings.

In practice, the exact ground state is not available. In ADC, it is expanded order by order in the fluctuation potential using M{\o}ller--Plesset perturbation theory. While the precise form is specified below, it matters here because, apart from the orbital vectors, the vectors of Eq.~\eqref{eq:si:vY} are orthonormal only if $\ket{\Psi_0}$ is the HF determinant. In general, expanding the states $\ket{\Psi^{N\pm1}_k}$ in this basis leads to the generalized eigenproblem
\begin{equation}
\mathcal{M}\mathbf{x} = \omega\,\mathcal{S}\mathbf{x} \;,
\label{eq:si:generalized_eigen}
\end{equation}
with the overlap matrix elements
\begin{equation}
\mathcal{S}_{IJ}=\braket{v(Y_I)|v(Y_J)}=\big\langle\{Y_I^\dagger,Y_J\}\big\rangle
\label{eq:si:S_IJ}
\end{equation}
and the Hamiltonian matrix elements
\begin{equation}
\mathcal{M}_{IJ}=\braket{v(Y_I)|\mathcal{H}|v(Y_J)}=\big\langle\{Y_I^\dagger,[H,Y_J]\}\big\rangle \;,
\label{eq:si:M_IJ}
\end{equation}
where $\langle\cdots\rangle$ denotes the expectation value in $\ket{\Psi_0}$. 

\subsubsection{Intermediate States in the Union Space}
This is a standard generalized eigenproblem, and there are many ways to convert this into a regular one. Here, we choose an orthonormal basis $\{\ket{\tilde{\Psi}_I}\}$ of the union space whose first vectors are the orbital vectors and whose remaining vectors are orthogonal to them. Such a basis can be constructed by orthonormalizing the initial basis class by class, by Gram--Schmidt orthogonalization to all lower classes followed by symmetric orthonormalization within each class. This gives the intermediate states of \citet{Schirmer_1991}. The matrix of $\mathcal{H}$ in the intermediate states, which we call $\tilde{\mathbf{M}}$, is the intermediate-state representation (ISR), and has matrix elements
\begin{equation}
\tilde{M}_{IJ} = \braket{\tilde{\Psi}_I|\mathcal{H}|\tilde{\Psi}_J} \;.
\label{eq:si:Mtilde}
\end{equation}

In the union space, $\tilde{\mathbf{M}}$ can be ordered by classes. The orbital block of its resolvent gives
\begin{equation}
\mathbf{G}(\omega)=\big[(\omega-\tilde{\mathbf{M}})^{-1}\big]_{00}=\big(\omega-\pmb\epsilon-\pmb\Sigma(\infty)-\pmb\Sigma_{\mathrm{dyn}}(\omega)\big)^{-1} \;.
\label{eq:si:resolvent}
\end{equation}
$\pmb{\epsilon}$ the diagonal matrix of orbital energies. The self-energy splits into a dynamic part $\pmb{\Sigma}_{\mathrm{dyn}}$ and a static part. The former vanishes at $\omega \to \infty$. The static part is the change of the mean field caused by the correlated density,\cite{Schirmer_1983,Schirmer_1998}
\begin{equation}
\Sigma_{pq}(\infty)=\sum_{rs}\langle pr||qs\rangle\,
\big(\gamma_{sr}-\gamma^{\mathrm{HF}}_{sr}\big) \;,
\label{eq:si:sinf}
\end{equation}
where $\gamma$ is the exact one-particle density matrix and $\gamma^{\mathrm{HF}}$ that of the reference.

\subsubsection{Dyson-ADC}
With the definition
\begin{equation}
\mathbf{G}^{(0)}(\omega) = \big(\omega-\pmb\epsilon\big)^{-1} \;,
\end{equation}
Eq.~\eqref{eq:si:resolvent} takes the form of a Dyson equation,
\begin{equation}
\mathbf{G}^{-1}(\omega) = \big[\mathbf{G}^{(0)}(\omega)\big]^{-1} -\pmb\Sigma(\infty) -  \pmb\Sigma_{\mathrm{dyn}}(\omega) \;.
\end{equation}
Its solution is the interacting propagator of the main text,
\begin{align}
G_{pq}(\omega) ={}& \sum_{n} \frac{\langle \Psi_0 | a_p | \Psi_n^{N+1} \rangle \langle \Psi_n^{N+1} | a_q^\dagger | \Psi_0 \rangle}{\omega - (E_n^{N+1} - E_0) + i\eta} \nonumber \\
&+ \sum_{k} \frac{\langle \Psi_0 | a_q^\dagger | \Psi_k^{N-1} \rangle \langle \Psi_k^{N-1} | a_p | \Psi_0 \rangle}{\omega - (E_0 - E_k^{N-1}) - i\eta} \;.
\label{eq:si:Lehmann}
\end{align}
The resolvent of the Hermitian matrix $\tilde{\mathbf{M}}$ gives the poles and residues of Eq.~\eqref{eq:si:Lehmann}; the infinitesimals $\pm i\eta$ of the time-ordered propagator are added afterwards, with the sign fixed by the sector of each pole. It is not obvious why the effective Hamiltonian constructed in this way, together with the Dyson equation, should generate Eq.~\eqref{eq:si:Lehmann}. The following conditions must hold, the first two being relatively obvious and the last one more subtle:
\begin{itemize}
\item The orbital vectors span all orbitals.
\item The fold contains the poles of both sectors.
\item The orbital vectors are orthonormal. In the union space their metric is the anticommutator, and orthonormality is simply the fermionic anticommutation relation,
\begin{equation}
\braket{v(a_p^\dagger)|v(a_q^\dagger)}=\langle\Psi_0|\{a_p,a_q^\dagger\}|\Psi_0\rangle=\delta_{pq} \;.
\end{equation}
\end{itemize}
Because this construction leads to a Dyson equation, it is referred to as Dyson ADC, to distinguish it from non-Dyson ADC.

\subsection{Non-Dyson ADC}
The non-Dyson form is the standard way to construct ADCs. It follows the same steps as above, but instead of working in the union space it constructs an effective Hamiltonian for the $(N+1)$- and the $(N-1)$-electron sectors separately, with the overlap and Hamiltonian matrix elements of Eqs.~\eqref{eq:si:S_IJ} and \eqref{eq:si:M_IJ} taken in one of the two spaces alone.

As the name suggests, non-Dyson ADC does not lead to a Dyson equation. One can still compute the orbital block of its resolvent as in Eq.~\eqref{eq:si:resolvent}, but the result is not of the form of Eq.~\eqref{eq:si:Lehmann}, because all three conditions above are violated:
\begin{itemize}
\item The orbital vectors do not span all orbitals. At zeroth order, they span only the occupied orbitals in the removal sector and only the virtual orbitals in the attachment sector.
\item The fold contains the poles of one sector only.
\item The orbital vectors are not orthonormal. In the removal space alone, for instance, their metric is only one of the two terms of the anticommutator,
\begin{equation}
\langle\Psi_0|a_p^\dagger a_q|\Psi_0\rangle=\gamma_{qp} \;,
\end{equation}
the one-particle density matrix. After orthonormalization, the spectroscopic amplitudes are no longer the unit matrix on the first class.
\end{itemize}

\subsection{Outline of this section}
The goal of this part of the Supporting Information is to discuss some subtle aspects of the Dyson construction of ADC. All of them have been discussed in the literature, most notably by \citet{Schirmer_1983}, \citet{tarantelli1992configuration} and \citet{Mertins_1996b}, but, to the best of our knowledge, they have never been reviewed together in one place.

While we have just described the intermediate-state construction of ADC in the union space, the original form of Dyson ADC derived by \citet{Schirmer_1983} was quite different. They started from the order-by-order expansion of the self-energy, the GF($n$) series, and enforced the correct spectral representation of the exact self-energy at every order of perturbation theory. The procedure is straightforward through third order but becomes subtle from fourth order on.

In the following, we derive two statements of the main text: that diagrammatic (Dyson) ADC($n$) is exactly the decoupled form of the union-space ISR of \citet{tarantelli1992configuration}, and that the two-sector ISR, which omits the cross block $\mathbf{X}$, agrees with diagrammatic ADC through third order but not beyond. We proceed as follows:
\begin{itemize}
\item We define ADC($n$) in Sec.~\ref{sec:si:setting}.
\item We describe the diagrammatic construction (Sec.~\ref{sec:si:diag}) and the ISR in the union space (Sec.~\ref{sec:si:isr}).
\item We compute the cross block and introduce the two-sector ISR in Sec.~\ref{sec:si:cross}.
\item We construct the decoupling with ADC(4) as the worked example (Sec.~\ref{sec:si:tc}).
\item We give the proof of the decoupling for general $n$ (Sec.~\ref{sec:si:proof}).
\end{itemize}

%-----------------------------------------------------------------------------
\subsection{\label{sec:si:setting}General definitions}
\subsubsection{Order counting}
For the following discussion we need the definition of order counting. We count orderes in the usual M{\o}ller--Plesset partitioning of the Hamiltonian. To do so, we scale the fluctuation potential by a coupling constant,
\begin{equation}
H(\lambda)=F+\lambda W \;,\qquad W=H-F \;,
\notag
\end{equation}
with $F$ the Fock operator. Every energy denominator of perturbation theory then becomes difference of orbital-energy sums and does not depend on $\lambda$. A
term built from $m$ matrix elements of $W$ carries a factor of $\lambda^m$. We call the Taylor coefficient of $\lambda^m$ the $m$th order and mark it by a superscript $(m)$. 

\subsubsection{Configurations and classes}
A configuration of class $k$ in the attachment sector has $k+1$ particles and $k$ holes relative to $\ket{\Phi_0}$. We denote it by $\nu$ and its zeroth-order energy, the sum of its particle energies minus the sum of its
hole energies, by $K^+_\nu$.

A configuration of class $k$ in the removal sector
has $k+1$ holes and $k$ particles. We denote it by $\mu$ and the energy needed to remove an electron into it, the sum of its particle energies minus the sum
of its hole energies, by $K^-_\mu$; the corresponding pole of the removal branch lies at minus $K^-_\mu$.

Class 0 holds the one-particle states, class 1 the 2p1h and 2h1p configurations, class 2 the 3p2h and 3h2p configurations, and so on. Combining a removal configuration of class $k$ with an attachment configuration of class $l$ gives a determinant of the neutral system with
$k+l+1$ particles and $k+l+1$ holes, so that
\begin{equation}
K^+_\nu+K^-_\mu=E^{(0)}_{(k+l+1)\mathrm{p}(k+l+1)\mathrm{h}}-E^{(0)}_0 \;.
\label{eq:si:gap}
\end{equation}
As an example, for two class-1 configurations this is a 3p3h excitation energy. It is positive and bounded away from zero by the HOMO-LUMO gap. In particular, at zeroth order every attachment pole lies above every removal pole, so the two sets of poles, which we call $P^+$ and $P^-$, are disjoint.

All constructions below deliver a static matrix and a fold
\begin{equation}
\mathbf{f}(\omega)=\mathbf{U}\,(\omega-\mathbf{K}-\mathbf{C})^{-1}\,\mathbf{U}^\dagger \;,
\label{eq:si:fold}
\end{equation}
where $\mathbf{K}$ is the diagonal matrix of zeroth-order pole positions and the Hermitian, frequency-independent $\mathbf{C}$ and the couplings $\mathbf{U}$ are power series in $\lambda$ that start at first order. Since $\mathbf{K}$ does not depend on $\lambda$, expanding the resolvent in powers of $\mathbf{C}$ (the identity below Eq.~\eqref{eq:GFn} of the main text) writes
each Taylor coefficient of the fold as a finite sum of chains
\begin{equation}
\mathbf{U}^{(a)}\,\mathbf{g}\,\mathbf{C}^{(b)}\,\mathbf{g}\cdots\mathbf{g}\,\mathbf{U}^{(c)\dagger} \;,
\qquad \mathbf{g}=(\omega-\mathbf{K})^{-1} \;,
\label{eq:si:chain}
\end{equation}
whose order is the sum of the orders of its links. Each chain is a rational function of $\omega$ that vanishes at infinity and has poles only at the diagonal entries of $\mathbf{K}$.

\subsubsection{Definition of ADC($n$)}
We define an ADC($n$) by the property that its static matrix and its fold reproduce the Taylor coefficients of the exact static and dynamic self-energy through
$\lambda^n$ at every frequency. 

\subsubsection{Split-Lemma}

The exact self-energy has the following important property. Any rational function $f(\omega)$ with the property $f(\omega) \rightarrow 0$ for $\omega \rightarrow \infty$ and has its poles in the union of $P^+$ and $P^-$ is the sum of exactly one function with poles in $P^+$ and one with poles in $P^-$, both vanishing at infinity. 
It is easy to see why: Suppose the split is not unique. Then, the function can be written as 
\begin{equation}
f(\omega) = f^+(\omega) + f^-(\omega) = g^+(\omega) + g^-(\omega) \;,
\end{equation}
with $f^+$, $g^+$ having poles in $P^+$ and $f^-$, $g^-$ having poles in $P^-$. Then, it follows that 
\begin{equation}
  f^+(\omega) - g^+(\omega) = g^-(\omega) - f^-(\omega) \;, 
\end{equation}
so the difference of two such splittings would have poles in both disjoint sets. This is only possible if it does not have any poles at all. The Taylor coefficients of the exact dynamic self-energy have their poles in $P^+$ and $P^-$, due to time-ordering. Therefore, for every order the dynamical self-energy can be written as 
\begin{equation}
\pmb{\Sigma}^{(m)}_{\mathrm{dyn}}=\pmb{\Sigma}^{(m)}_{+}+\pmb{\Sigma}^{(m)}_{-} \;.
\label{eq:si:pmsplit}
\end{equation}
That the self-energy can be written this way also follows from its Goldstone expansion, that shows that each time ordering of a self-energy diagram uniquely belongs to one sector. If the electron enters at the lower
external vertex, the system holds $N+1$ electrons between the external vertices. Therefore, every frequency-dependent energy denominator is that of an attachment configuration. In the reverse time order it is that of a removal configuration. 

\subsection{\label{sec:si:diag}Diagrammatic (Dyson) ADC}
\subsubsection{General Procedure}
Schirmer, Cederbaum and Walter write each part of the dynamic self-energy as a fold over its own configuration space,\cite{Schirmer_1983}
\begin{align}
\pmb{\Sigma}_{+}(\omega)&=\mathbf{U}^{+}\big(\omega-\mathbf{K}^{+}-\mathbf{C}^{+}\big)^{-1}\mathbf{U}^{+\dagger} \;,
\notag\\
\pmb{\Sigma}_{-}(\omega)&=\mathbf{U}^{-}\big(\omega+\mathbf{K}^{-}+\mathbf{C}^{-}\big)^{-1}\mathbf{U}^{-\dagger} \;,
\notag
\end{align}
the two branches of Eq.~\eqref{eq:SigmaADC} of the main text, over the 2p1h, 3p2h, \dots\ and the 2h1p, 3h2p, \dots\ configurations. Expanding in $\lambda$ gives, for the attachment part with the sector label dropped,
\begin{align}
\pmb{\Sigma}^{(2)}={}&\mathbf{U}^{(1)}\mathbf{g}\,\mathbf{U}^{(1)\dagger} \;,
  \label{eq:adc2}\\
\pmb{\Sigma}^{(3)}={}&\mathbf{U}^{(1)}\mathbf{g}\,\mathbf{U}^{(2)\dagger}+\mathbf{U}^{(2)}\mathbf{g}\,\mathbf{U}^{(1)\dagger}
+\mathbf{U}^{(1)}\mathbf{g}\,\mathbf{C}^{(1)}\mathbf{g}\,\mathbf{U}^{(1)\dagger} \;,
  \label{eq:adc3}\\
\pmb{\Sigma}^{(4)}={}&\mathbf{U}^{(1)}\mathbf{g}\,\mathbf{U}^{(3)\dagger}+\mathbf{U}^{(3)}\mathbf{g}\,\mathbf{U}^{(1)\dagger}
+\mathbf{U}^{(2)}\mathbf{g}\,\mathbf{U}^{(2)\dagger}
\notag\\
&+\mathbf{U}^{(1)}\mathbf{g}\,\mathbf{C}^{(1)}\mathbf{g}\,\mathbf{U}^{(2)\dagger}
+\mathbf{U}^{(2)}\mathbf{g}\,\mathbf{C}^{(1)}\mathbf{g}\,\mathbf{U}^{(1)\dagger}
\notag\\
&+\mathbf{U}^{(1)}\mathbf{g}\,\mathbf{C}^{(2)}\mathbf{g}\,\mathbf{U}^{(1)\dagger}
+\mathbf{U}^{(1)}\mathbf{g}\,\mathbf{C}^{(1)}\mathbf{g}\,\mathbf{C}^{(1)}\mathbf{g}\,\mathbf{U}^{(1)\dagger} \;.
  \label{eq:adc4}
\end{align}
This is Eq.~\eqref{eq:GFn} of the main text completed at fourth order. ADC($n$) follows by requiring that these coefficients equal the exact attachment and removal parts of Eq.~\eqref{eq:si:pmsplit} through order $n$. In these matching equations the new unknowns at order $m+1$ are $\mathbf{U}^{(m)}$ and $\mathbf{C}^{(m-1)}$, and they enter linearly, so the equations are solved order by order. The construction needs two additional rules. First, ADC($n$) keeps the classes up to $\lfloor n/2\rfloor$. Second, the block between classes $k$ and $l$ is carried through order $n-k-l$. These rules have been given by \citet{Schirmer_1983}, but they can be derived in the intermediate state representation. This derivation is postponed to Section~\ref{sec:si:isr}. The exactness and uniqueness of the split in Eq.~\eqref{eq:si:pmsplit} allow for the separate matching of the addition and removal branches of the dynamical self-energy.

\subsubsection{Second order}
At second order, comparing with Eq.~\eqref{eq:Sigma2} of the main text gives the first-order couplings
\begin{equation}
U^{+(1)}_{p,(abk)}=\langle pk||ab\rangle \;,\qquad
U^{-(1)}_{p,(ija)}=\langle pa||ij\rangle \;.
\notag
\end{equation}
This corresponds to Goldstone diagrams that have only a single, frequency-dependent denominator. At second order, this is the only denominator form.

\subsubsection{Third order}
At third order, the Goldstone expansion of the self-energy produces three classes of diagrams. Each of them has two denominators. The first class has two static denominators. Those are then part of the static self-energy. The second class has one static and one dynamic denominator. These terms then belong to the first two terms of the \emph{r.h.s.} of Eq.~\eqref{eq:adc3}. The static denominators are absorbed into the couplings and combine with a vertex into a first-order doubles amplitude,
\begin{equation}
t^{ab}_{ij}=\frac{\langle ab||ij\rangle}{\epsilon_i+\epsilon_j-\epsilon_a-\epsilon_b} \;,
\label{eq:si:t2}
\end{equation}
Finally, the third class of diagrams has two dynamical denominators. These enter the third term of the \emph{r.h.s.} of Eq.~\eqref{eq:adc3}.

\subsubsection{Gauge Freedom}
The matrices are determined by the matching only up to a unitary transformation that reflects gauge freedom.\cite{tarantelli1992configuration} Take any real antisymmetric matrix $\mathbf{T}$ of first or higher order that acts within one sector,
\begin{equation}
  \begin{aligned}
\pmb{\Sigma}_{+}(\omega) = & \mathbf{U}^{+}\big(\omega-\mathbf{K}^{+}-\mathbf{C}^{+}\big)^{-1}\mathbf{U}^{+\dagger} \\
= & 
\left(
\mathbf{U}^{+}e^{\mathbf{T}}\right) 
\left(
e^{-\mathbf{T}}\big(\omega-\mathbf{K}^{+}-\mathbf{C}^{+}\big)^{-1}
e^{\mathbf{T}}
\right)
\left(
e^{-\mathbf{T}}
\mathbf{U}^{+\dagger}
\right) \;,
  \end{aligned}
\label{eq:si:gauge}
\end{equation}
and analogous for the $-$-branch. Physically, this reflects the fact that a chain with two different resolvents decomposes into two single poles, but their residues cannot tell whether they came from $\mathbf{C}$ or from $\mathbf{U}$. Importantly, since $\mathbf{T}$ is at least of first order,
\begin{equation}
  e^{-\mathbf{T}}(\mathbf{K}+\mathbf{C})e^{\mathbf{T}} = \mathbf{K}+\mathbf{C}' \;,
\label{eq:si:gaugeK}
\end{equation}
where $\mathbf{C}'$ collects $\mathbf{C}$ and all commutators with $\mathbf{T}$, which are of first or higher order. Hence, this transformation only defines alternative couplings $\mathbf{U}$ and class interaction blocks $\mathbf{C}$, but no new zeroth-order poles. The diagrammatic construction implicitly fixes this gauge by a structural rule. A time ordering with $r$ frequency-dependent denominators is assigned to the chain with $r$ resolvents, as explained above for third order.

\subsubsection{Fourth Order}
At fourth order the matching is less straightforward. The reason is that starting from fourth order diagrams with frequency-dependent denominators at a 4p3h or 4h3p configuration show up. This seems to imply that ADC(4) must contain this configuration space, i.e. going from ADC(3) to ADC(4) would imply that the configuration space is extended by two classes. This seems odd, and can indeed be avoided. As shown in Ref.~\citenum{Schirmer_1983}, these diagrams always come in groups of four, whose residues at those poles cancel. Each group then sums to terms with poles at 2p1h/2h1p configurations only. The fourth-order terms with the weight of a 3p3h times a 2p2h denominator come from a separate class of time orderings, those with a 3p3h configuration in flight between two class-1 configurations, and they carry the content of the cross block of Sec.~\ref{sec:si:cross}. Beyond fourth order, this matching becomes cumbersome, even though it can be automated. The ISR overcomes the laborious matching task beyond third order.\cite{Schirmer_1991} Section~\ref{sec:si:cross} shows that its two-sector form avoids it by leaving out exactly the content with the 3p3h weight.

\subsection{\label{sec:si:isr}The union space effective hamiltonian}

\subsubsection{General Properties}
Written out block by block, the ISR of the matrix $\tilde{\mathbf{M}}$ with the matrix elements of Eq.~\eqref{eq:si:Mtilde} becomes
\begin{equation}
\tilde{\mathbf{M}}=
\begin{pmatrix}
\pmb{\epsilon}+\pmb{\Sigma}(\infty) & \mathbf{U}^{-} & \mathbf{U}^{+}\\
\mathbf{U}^{-\dagger} & -\mathbf{K}^{-}-\mathbf{C}^{-} & \mathbf{X}\\
\mathbf{U}^{+\dagger} & \mathbf{X}^{\dagger} & \mathbf{K}^{+}+\mathbf{C}^{+}
\end{pmatrix} ,
\label{eq:si:union}
\end{equation}
This is Eq.~\eqref{eq:ISRunion} of the main text.
Downfolding of this expression to the orbital block and comparison with Eq.~\eqref{eq:si:union} gives a dynamical self-energy of the form
\begin{equation}
\pmb{\Sigma}_{\mathrm{dyn}}(\omega)=\tilde{\mathbf{U}}\,(\omega-\tilde{\mathbf{M}}_{\mathrm{c}})^{-1}\,\tilde{\mathbf{U}}^\dagger \;,
\label{eq:si:unionexact}
\end{equation}
where $\tilde{\mathbf{U}}$ collects both couplings and
\begin{equation}
\tilde{\mathbf{M}}_{\mathrm{c}} = 
\begin{pmatrix}
 -\mathbf{K}^{-}-\mathbf{C}^{-} & \mathbf{X}\\
\mathbf{X}^{\dagger} & \mathbf{K}^{+}+\mathbf{C}^{+}
\end{pmatrix} \;.
\end{equation}
Due  the presence of the block $\mathbf{X}$ it is not yet of the form of Eq.~\eqref{eq:si:pmsplit}, i.e. it does not split into electron addition and removal branch. In other words, it does not havew the form of a physical self-energy.

To bring it into this form, we need an orthogonal transformation that leaves the orbital vectors unchanged and mixes only the configurations. Such a transformation looks like
\begin{equation}
\mathbf{V}=
\begin{pmatrix}
\mathbf{1} & \mathbf{0} & \mathbf{0}\\
\mathbf{0} & \mathbf{V}^{--} & \mathbf{V}^{-+}\\
\mathbf{0} & \mathbf{V}^{+-} & \mathbf{V}^{++}
\end{pmatrix} ,
\label{eq:si:Vform}
\end{equation}
with blocks chosen such that the transformed matrix has no cross block,
\begin{equation}
\mathbf{V}^\dagger\tilde{\mathbf{M}}\mathbf{V}=
\begin{pmatrix}
\pmb{\epsilon}+\pmb{\Sigma}(\infty) & \mathbf{U}^{-\prime} & \mathbf{U}^{+\prime}\\
\mathbf{U}^{-\prime\dagger} & -\mathbf{K}^{-}-\mathbf{C}^{-\prime} & \mathbf{0}\\
\mathbf{U}^{+\prime\dagger} & \mathbf{0} & \mathbf{K}^{+}+\mathbf{C}^{+\prime}
\end{pmatrix} .
\label{eq:si:Vdecoupled}
\end{equation}
Since $\mathbf{V}$ does not touch the orbital vectors, it leaves the orbital block of the resolvent in Eq.~\eqref{eq:si:Gblock} and hence the self-energy
unchanged. Block inversion of the transformed matrix now gives two separate
folds,
\begin{align}
\pmb{\Sigma}_{\mathrm{dyn}}(\omega)={}&
\mathbf{U}^{+\prime}\big(\omega-\mathbf{K}^{+}-\mathbf{C}^{+\prime}\big)^{-1}\mathbf{U}^{+\prime\dagger}
\notag\\
&+\mathbf{U}^{-\prime}\big(\omega+\mathbf{K}^{-}+\mathbf{C}^{-\prime}\big)^{-1}\mathbf{U}^{-\prime\dagger} \;,
\label{eq:si:Vfolds}
\end{align}
which is the form of Eq.~\eqref{eq:SigmaADC} of the main text. The blocks of $\mathbf{V}$ are specified in Sec.~\ref{sec:si:tc}.

The transformation $\mathbf{V}$ is of the same kind as the gauge transformation of Eqs.~\eqref{eq:si:gauge} and~\eqref{eq:si:gaugeK}. Both leave the orbitals fixed and change only the couplings and the configuration block, so that the fold is unchanged, and both reduce to the identity at zeroth order, so that no new zeroth-order poles appear. They differ in which configurations they mix. The gauge transformation acts within one sector, whereas $\mathbf{V}$ must also mix the two sectors through $\mathbf{V}^{-+}$ and $\mathbf{V}^{+-}$ to remove $\mathbf{X}$. The decoupling condition of Eq.~\eqref{eq:si:Vdecoupled} fixes only this mixing. Any further rotation within each sector, i.e., a gauge transformation of Eq.~\eqref{eq:si:gauge},
preserves the decoupled form, so the blocks $\mathbf{V}^{--}$ and $\mathbf{V}^{++}$ are not determined by the decoupling. This is the same freedom that the diagrammatic construction fixes by its structural rule.

\subsubsection{Order Truncation}

To truncate this matrix consistently, we need to know at which order in $\lambda$ each of its blocks starts. This tells us which classes and which
orders of each block ADC($n$) must keep, and it derives the order rule quoted in Sec.~\ref{sec:si:diag}. Both the decoupling of Sec.~\ref{sec:si:tc} and the proof of Sec.~\ref{sec:si:proof} rely on these orders.

The matrix elements become power series in $\lambda$ when the Rayleigh--Schr\"odinger expansion in the M{\o}ller--Plesset partitioning, here in intermediate normalization,
\begin{equation}
\ket{\Psi_0} = \ket{\Phi_0} + \lambda\ket{\Psi^{(1)}} + \lambda^2\ket{\Psi^{(2)}} + \mathcal{O}(\lambda^3) \;,
\label{eq:si:psiexp}
\end{equation}
is inserted. Because this state is not normalized, the expectation values of Eq.~\eqref{eq:si:M_IJ} must be divided by its norm. Writing the same state in exponential form, $\ket{\Psi_0}=e^{\hat{T}}\ket{\Phi_0}$, with $\hat{T}$ the cluster operator of all ranks, gives
\begin{equation}
\mathcal{M}_{IJ}=\frac{1}{\braket{\Psi_0|\Psi_0}}\sum_{m,n}\frac{1}{m!\,n!}\braket{\Phi_0|(\hat{T}^\dagger)^m\,\{Y_I^\dagger,[H,Y_J]\}\,\hat{T}^n|\Phi_0} \;,
\label{eq:si:Mexp}
\end{equation}
where $m$ and $n$ count the powers of the cluster operator; each power carries all orders of its amplitudes. The question is now which orders enter each block between classes $k$ and $l$, for instance the orbital block $k=l=0$, the coupling of the orbitals to the 2p1h configurations, and so on. Let us first write out Eq.~\eqref{eq:si:psiexp} explicitly. Its first-order term is
\begin{equation}
\ket{\Psi^{(1)}} = \frac14\sum_{ijab} t^{ab}_{ij}\, a_a^\dagger a_b^\dagger a_j a_i \ket{\Phi_0} \;, 
\label{eq:si:psi1}
\end{equation}
It contains the first-order double excitations with the amplitudes of Eq.~\eqref{eq:si:t2}.
The second-order wave function contains singles, doubles, connected triples, and the disconnected quadruples of the first-order operator,
\begin{align}
\ket{\Psi^{(2)}} ={}& \Big[\sum_{ia} t^{a(2)}_{i}\, a_a^\dagger a_i
+ \frac14\sum_{ijab} t^{ab(2)}_{ij}\, a_a^\dagger a_b^\dagger a_j a_i
\notag\\
&+ \frac{1}{36}\sum_{ijkabc} t^{abc(2)}_{ijk}\, a_a^\dagger a_b^\dagger a_c^\dagger a_k a_j a_i
+ \frac12\big(\hat{T}^{(1)}\big)^2\Big]\ket{\Phi_0} \;,
\label{eq:si:psi2}
\end{align}
where $\hat{T}^{(1)}$ is the double-excitation operator of Eq.~\eqref{eq:si:psi1}, which creates $\ket{\Psi^{(1)}}$ from $\ket{\Phi_0}$. The second-order amplitudes are
\begin{align}
D^{a}_{i}\, t^{a(2)}_{i} ={}& \frac12\sum_{jbc}\langle aj||bc\rangle\, t^{bc}_{ij}
- \frac12\sum_{jkb}\langle jk||ib\rangle\, t^{ab}_{jk} \;,
\label{eq:si:t1two}\\
D^{ab}_{ij}\, t^{ab(2)}_{ij} ={}& \frac12\sum_{cd}\langle ab||cd\rangle\, t^{cd}_{ij}
+ \frac12\sum_{kl}\langle kl||ij\rangle\, t^{ab}_{kl}
\notag\\
&+ P(ij)P(ab)\sum_{kc}\langle kb||cj\rangle\, t^{ac}_{ik} \;,
\label{eq:si:t2two}\\
D^{abc}_{ijk}\, t^{abc(2)}_{ijk} ={}& P(i/jk)P(a/bc)\Big[\sum_{d}\langle di||bc\rangle\, t^{ad}_{jk}
\notag\\
&- \sum_{l}\langle la||jk\rangle\, t^{bc}_{il}\Big] \;,
\label{eq:si:t3two}
\end{align}
with the orbital-energy differences
\begin{align}
D^{a}_{i} &= \epsilon_i - \epsilon_a \;,
\notag\\
D^{ab}_{ij} &= \epsilon_i + \epsilon_j - \epsilon_a - \epsilon_b \;,
\notag\\
D^{abc}_{ijk} &= \epsilon_i + \epsilon_j + \epsilon_k - \epsilon_a - \epsilon_b - \epsilon_c \;,
\notag
\end{align}
and the antisymmetrizers
\begin{align}
P(ij)\,f(i,j) &= f(i,j) - f(j,i) \;,
\notag\\
P(i/jk)\,f(i,j,k) &= f(i,j,k) - f(j,i,k) - f(k,j,i) \;.
\notag
\end{align}

\subsubsection{Oder truncation Rules}
We can now simply work out each block using Wick's theorem, and find thus an "empirical" order rule by counting the orders that survive in each block. It is however also possible to derive a general order rule first, which provides us with a lower bound for the order at which each block starts:\cite{Schirmer_1991,Mertins_1996}
\begin{itemize}
\item 
The coupling of the orbitals to class $k$ starts at order $k$
\item The block between classes $k$ and $l$ starts at order $|k-l|$. 
\item Within a class the diagonal $\mathbf{K}$ is of zeroth order and $\mathbf{C}$ starts at first order. 
\end{itemize}
The cross blocks between the two sectors do not follow this rule in general. Their orders will be derived in Sec.~\ref{sec:si:cross}.

It is also important to know until which order of perturbation theory we need each block. This is illustrated with an example. Consider the attachment block of Eq.~\eqref{eq:si:Vform}, i.e. after the removal and attachment sectors have been decoupled, restricted to classes up to 2,
\begin{equation}
\begin{pmatrix}
\pmb{\epsilon}+\pmb{\Sigma}(\infty) & \mathbf{U}^{+}_1 & \mathbf{U}^{+}_2\\
\mathbf{U}^{+\dagger}_1 & \mathbf{K}^{+}_1+\mathbf{C}^{+}_{11} & \mathbf{C}^{+}_{12}\\
\mathbf{U}^{+\dagger}_2 & \mathbf{C}^{+}_{21} & \mathbf{K}^{+}_2+\mathbf{C}^{+}_{22}
\end{pmatrix} ,
\label{eq:si:Mclasses}
\end{equation}
with the following starting order of each block due to the canonical order relations,
\begin{equation}
\text{starting orders} = \quad
\begin{pmatrix}
0 & 1 & 2\\
1 & 0\,(\mathbf{K}),\,1\,(\mathbf{C}) & 1\\
2 & 1 & 0\,(\mathbf{K}),\,1\,(\mathbf{C})
\end{pmatrix} \;.
\end{equation}
To find the maximum order at which a block can contribute to ADC($n$), we downfold Eq.~\eqref{eq:si:Mclasses} one class at a time onto the lower classes. Folding class 2 onto the orbitals and class 1 gives
\begin{equation}
\begin{pmatrix}
\pmb{\epsilon}+\pmb{\Sigma}(\infty)+\mathbf{U}^{+}_2\mathbf{g}_2\mathbf{U}^{+\dagger}_2 & \mathbf{U}^{+}_1+\mathbf{U}^{+}_2\mathbf{g}_2\mathbf{C}^{+}_{21}\\
\mathbf{U}^{+\dagger}_1+\mathbf{C}^{+}_{12}\mathbf{g}_2\mathbf{U}^{+\dagger}_2 & \mathbf{K}^{+}_1+\mathbf{C}^{+}_{11}+\mathbf{C}^{+}_{12}\mathbf{g}_2\mathbf{C}^{+}_{21}
\end{pmatrix} ,
\label{eq:si:fold2}
\end{equation}
with $\mathbf{g}_2=(\omega-\mathbf{K}^{+}_2-\mathbf{C}^{+}_{22})^{-1}$. Folding class 1 onto the orbitals then gives the attachment branch of the self-energy,
\begin{equation}
\begin{aligned}
\pmb{\Sigma}^{+}_{\mathrm{dyn}}(\omega)= & \big(\mathbf{U}^{+}_1+\mathbf{U}^{+}_2\mathbf{g}_2\mathbf{C}^{+}_{21}\big)\,\mathbf{g}_1\,\big(\mathbf{U}^{+}_1+\mathbf{U}^{+}_2\mathbf{g}_2\mathbf{C}^{+}_{21}\big)^\dagger \\
& +\mathbf{U}^{+}_2\,\mathbf{g}_2\,\mathbf{U}^{+\dagger}_2 \;,
\end{aligned}
\label{eq:si:fold1}
\end{equation}
with 
\begin{equation}
\mathbf{g}_1=(\omega-\mathbf{K}^{+}_1-\mathbf{C}^{+}_{11}-\mathbf{C}^{+}_{12}\mathbf{g}_2\mathbf{C}^{+}_{21})^{-1} \;.
\end{equation}
We can now deduce from this example up to which order each block is needed:
\begin{itemize}
\item The static self-energy is not changed by any fold, so ADC($n$) needs it through order $n$.
\item The coupling $\mathbf{U}_k$ enters the self-energy only through products $\mathbf{U}_k\mathbf{g}_k\mathbf{U}_k^\dagger$. A term of order $q$ in $\mathbf{U}_k$ is multiplied there with $\mathbf{U}_k^\dagger$, which starts at order $k$, and with $\mathbf{g}_k$, which starts at zeroth order, so the product is of order at least $q+k \geq 2k$.

\item The class block $\mathbf{C}_{kk}$ enters inside $\mathbf{g}_k$, between two couplings $\mathbf{U}_k$ that start at order $k$. A term of order $p$ in $\mathbf{C}_{kk}$ contributes at order $2k+p$. Therefore, $\mathbf{C}_{kk}$ is needed only through order $n-2k$. Class $k$ is needed at all only if its zeroth-order diagonal contributes, $2k\le n$.
\item The class-class coupling $\mathbf{C}_{kl}$, $k<l$, enters when class $l$ is folded into class $k$, both in the corrected coupling $\mathbf{U}_k+\mathbf{U}_l\mathbf{g}_l\mathbf{C}_{lk}$ and in $\mathbf{g}_k$ through $\mathbf{C}_{kl}\mathbf{g}_l\mathbf{C}_{lk}$. In either place, the term that contains $\mathbf{C}_{kl}$ at order $p$ must reach class $k$ from the orbitals and return from class $l$, which costs at least $k$ and $l$ orders, so it contributes at order $k+p+l$. Therefore, $\mathbf{C}_{kl}$ is needed through order $n-k-l$.
\end{itemize}

The result, together with its generalization to ADC($n$) is given in Table~\ref{tab:si:sameblocks}.
\begin{table}[htb]
\caption{Same-type blocks of ADC($n$) under the order rule and the orders carried. Class 1: 2p1h or 2h1p; class 2: 3p2h or 3h2p. A block is absent when its upper order is below its lower one. The orbital block is carried through order $n$.}
\label{tab:si:sameblocks}
\begin{tabular}{lccccc}
\hline\hline
block & ADC(2) & ADC(3) & ADC(4) & ADC(5) & ADC($n$)\\
\hline
$\mathbf{U}$, class 1 & 1 & 1--2 & 1--3 & 1--4 & $1$--$(n-1)$\\
$\mathbf{K}+\mathbf{C}$, class 1 & 0 & 0--1 & 0--2 & 0--3 & $0$--$(n-2)$\\
$\mathbf{U}$, class 2 & -- & -- & 2 & 2--3 & $2$--$(n-2)$\\
$\mathbf{C}$, classes 1--2 & -- & -- & 1 & 1--2 & $1$--$(n-3)$\\
$\mathbf{K}+\mathbf{C}$, class 2 & -- & -- & 0 & 0--1 & $0$--$(n-4)$\\
\hline
$\mathbf{U}$, class $k$ & & & & & $k$--$(n-k)$\\
$\mathbf{C}$, classes $k$--$l$ & & & & & $|k-l|$--$(n-k-l)$\\
$\mathbf{K}+\mathbf{C}$, class $k$ & & & & & $0$--$(n-2k)$\\
\hline\hline
\end{tabular}
\end{table}

\subsection{Two sectors and the cross block}
\label{sec:si:cross}

This section explicitly demonstrates the claim of the main text that the cross couplings between attachment and removal classes vanish at zeroth and first order only, and are non-zero starting from second order.

The explicit expression for the couplings
elements is
\begin{align}
X_{\mu\nu}={}&\langle\Psi_0|\{O^-_\mu,[H,O^+_\nu]\}|\Psi_0\rangle
\notag\\
={}&\langle O^{-\dagger}_\mu\Psi_0|H-E_0|O^+_\nu\Psi_0\rangle
-\langle O^-_\mu\Psi_0|H-E_0|O^{+\dagger}_\nu\Psi_0\rangle \;.
\label{eq:si:Xraw}
\end{align}
We now compute $\mathbf{X}$ order by order for a class-1 pair with removal
string
\begin{equation}
  \begin{aligned}
O^-_\mu = & a_a^\dagger a_j a_i \\
O^+_\nu = & a_c^\dagger a_d^\dagger a_k  \;.
  \end{aligned}
\label{eq:si:Xpair}
\end{equation}
We also need the shifted Hamiltonian,
\begin{equation}
H - E_0 = \big(F - E^{(0)}_0\big) + \lambda\big(W - E^{(1)}_0\big) + \mathcal{O}(\lambda^2) \;.
\notag
\end{equation}
We denote the determinants reached from $\ket{\Phi_0}$ by the strings as
\begin{equation}
\begin{aligned}
\ket{\Phi^{cd}_{k}} &= a_c^\dagger a_d^\dagger a_k\ket{\Phi_0} \;, &
\ket{\Phi^{a}_{ij}} &= a_a^\dagger a_j a_i\ket{\Phi_0} \;, \\
\ket{\Phi^{b}} &= a_b^\dagger\ket{\Phi_0} \;, &
\ket{\Phi_{l}} &= a_l\ket{\Phi_0} \;.
\end{aligned}
\notag
\end{equation}
The Fock operator is diagonal in determinants,
\begin{equation}
\braket{\Phi_I|F|\Phi_J} = \delta_{IJ}\,E^{(0)}_I \;,
\label{eq:si:Fdiag}
\end{equation}
$E^{(0)}_I$ is the sum of the occupied orbital energies of $\ket{\Phi_I}$.

\subsubsection{Zeroth Order}

At zeroth order, the ground state is $\ket{\Phi_0}$, and the adjoint strings
in Eq.~\eqref{eq:si:Xraw} annihilate it, because both contain the annihilator
of an empty virtual orbital,
\begin{align}
O^{-\dagger}_\mu\ket{\Phi_0} &= a_i^\dagger a_j^\dagger a_a\ket{\Phi_0} = 0 \;,
\notag\\
O^{+\dagger}_\nu\ket{\Phi_0} &= a_k^\dagger a_d a_c\ket{\Phi_0} = 0 \;.
\label{eq:si:adj0}
\end{align}
Each of the two terms of Eq.~\eqref{eq:si:Xraw} contains one of these vectors, so
\begin{equation}
X^{(0)}_{\mu\nu} = 0 \;.
\notag
\end{equation}
In union-space language, with an HF ground state each intermediate state lives
in one sector only, and the composite Hamiltonian does not connect the two
sectors.

\subsubsection{First Order}

At first order, correlation gives each vector a component in the other sector.
The upper component of the removal-type vector becomes a one-particle state,
and the lower component of the attachment-type vector a one-hole state,
\begin{align}
O^{-\dagger}_\mu\ket{\Psi^{(1)}}&=\sum_b t^{ab}_{ij}\,\ket{\Phi^{b}} \;,
\label{eq:si:admix1}\\
O^{+\dagger}_\nu\ket{\Psi^{(1)}}&=-\sum_l t^{cd}_{kl}\,\ket{\Phi_{l}} \;.
\label{eq:si:admix2}
\end{align}
Collecting the first-order terms of Eq.~\eqref{eq:si:Xraw}, every term that
contains one of the vectors of Eq.~\eqref{eq:si:adj0} vanishes, and only the
admixtures combined with the zeroth-order Hamiltonian remain,
\begin{align}
X^{(1)}_{\mu\nu} ={}& \braket{O^{-\dagger}_\mu\Psi^{(1)}|F-E^{(0)}_0|O^+_\nu\Phi_0}
- \braket{O^-_\mu\Phi_0|F-E^{(0)}_0|O^{+\dagger}_\nu\Psi^{(1)}}
\notag\\
={}& \sum_b t^{ab}_{ij}\,\braket{\Phi^{b}|F-E^{(0)}_0|\Phi^{cd}_{k}}
+ \sum_l t^{cd}_{kl}\,\braket{\Phi^{a}_{ij}|F-E^{(0)}_0|\Phi_{l}} \;.
\notag
\end{align}
Both matrix elements connect different determinants, a one-particle determinant
with a 2p1h determinant in the first sum and a 2h1p determinant with a one-hole
determinant in the second. By Eq.~\eqref{eq:si:Fdiag} they vanish, so that\cite{tarantelli1992configuration}
\begin{equation}
X^{(1)}_{\mu\nu} = 0 \;,
\notag
\end{equation}

\subsubsection{Second Order}

At second order, we collect all terms of Eq.~\eqref{eq:si:Xraw} whose orders
in the bra, the Hamiltonian, and the ket add up to two,
\begin{align}
X^{(2)}_{\mu\nu} = \sum_{p+q+r=2} \Big[&
\braket{O^{-\dagger}_\mu\Psi^{(p)}|(H-E_0)^{(q)}|O^+_\nu\Psi^{(r)}}
\notag\\
&- \braket{O^-_\mu\Psi^{(p)}|(H-E_0)^{(q)}|O^{+\dagger}_\nu\Psi^{(r)}}\Big] \;,
\label{eq:si:X2sum}
\end{align}
with $\ket{\Psi^{(0)}}$ the HF determinant and the orders of the shifted
Hamiltonian
\begin{align}
(H-E_0)^{(0)} &= F-E^{(0)}_0 \;,
\notag\\
(H-E_0)^{(1)} &= W-E^{(1)}_0 \;,
\notag\\
(H-E_0)^{(2)} &= -E^{(2)}_0 \;.
\notag
\end{align}
By Eq.~\eqref{eq:si:adj0}, every term with $p$ equal to zero in the first part
and every term with $r$ equal to zero in the second part vanishes. Two of the
remaining terms have both vectors at first order. In the first, the bra is the
one-particle state of Eq.~\eqref{eq:si:admix1}, while the ket,
\begin{equation}
O^+_\nu\ket{\Psi^{(1)}} = \frac14\sum_{mnef} t^{ef}_{mn}\,
a_c^\dagger a_d^\dagger a_k\, a_e^\dagger a_f^\dagger a_n a_m\ket{\Phi_0} \;,
\notag
\end{equation}
is a combination of 4p3h determinants. In the second, the ket is the one-hole
state of Eq.~\eqref{eq:si:admix2}, while the bra $O^-_\mu\ket{\Psi^{(1)}}$ is a
combination of 4h3p determinants. The Fock operator is diagonal,
Eq.~\eqref{eq:si:Fdiag}, so both terms vanish,
\begin{equation}
\braket{O^{-\dagger}_\mu\Psi^{(1)}|F-E^{(0)}_0|O^+_\nu\Psi^{(1)}} =
\braket{O^-_\mu\Psi^{(1)}|F-E^{(0)}_0|O^{+\dagger}_\nu\Psi^{(1)}} = 0 \;.
\notag
\end{equation}
Four terms remain,
\begin{align}
X^{(2)}_{\mu\nu} ={}& \braket{O^{-\dagger}_\mu\Psi^{(2)}|F-E^{(0)}_0|O^+_\nu\Phi_0}
- \braket{O^-_\mu\Phi_0|F-E^{(0)}_0|O^{+\dagger}_\nu\Psi^{(2)}}
\notag\\
&+ \braket{O^{-\dagger}_\mu\Psi^{(1)}|W|O^+_\nu\Phi_0}
- \braket{O^-_\mu\Phi_0|W|O^{+\dagger}_\nu\Psi^{(1)}} \;.
\label{eq:si:X2terms}
\end{align}
In the first two terms, the Fock operator acts on the determinants
$\ket{\Phi^{cd}_k}$ and $\ket{\Phi^{a}_{ij}}$ and gives their zeroth-order
energies relative to the ground state, $K^+_\nu$ and $K^-_\mu$. What remains
is the overlap of $\ket{\Psi^{(2)}}$ with the triply excited determinant
obtained by applying both strings to the reference (holes $i,j,k$, particles
$a,c,d$),
\begin{equation}
\ket{T_{\mu\nu}} = O^-_\mu O^+_\nu\ket{\Phi_0} \;, \qquad
c^{(2)}_{\mu\nu} = \braket{T_{\mu\nu}|\Psi^{(2)}} \;,
\notag
\end{equation}
where we used that the two odd strings anticommute in the second term. In the
last two terms, the admixtures of Eqs.~\eqref{eq:si:admix1}
and~\eqref{eq:si:admix2} meet one matrix element of $W$ between determinants
that differ by two spin orbitals, which the Slater--Condon rules give as an
antisymmetrized integral. Together,
\begin{align}
X^{(2)}_{\mu\nu}={}&(K^+_\nu+K^-_\mu)\,c^{(2)}_{\mu\nu}
+\sum_b t^{ab}_{ij}\langle bk||cd\rangle
\notag\\
&+\sum_l t^{cd}_{kl}\langle al||ij\rangle \;.
\label{eq:si:X2}
\end{align}
The triples coefficient follows from the Rayleigh--Schr\"odinger recursion.
The zeroth-order energy of $\ket{T_{\mu\nu}}$ relative to the ground state is
the sum of $K^+_\nu$ and $K^-_\mu$, Eq.~\eqref{eq:si:gap}, so
\begin{equation}
c^{(2)}_{\mu\nu} = -\frac{\braket{T_{\mu\nu}|W|\Psi^{(1)}}}{K^+_\nu+K^-_\mu} \;,
\notag
\end{equation}
and the triples denominator cancels in Eq.~\eqref{eq:si:X2},
\begin{align}
X^{(2)}_{\mu\nu}={}&-\braket{T_{\mu\nu}|W|\Psi^{(1)}}
+\sum_b t^{ab}_{ij}\langle bk||cd\rangle
\notag\\
&+\sum_l t^{cd}_{kl}\langle al||ij\rangle \;.
\label{eq:si:X2final}
\end{align}
The whole cross block at
second order is a contraction of integrals with first-order doubles
amplitudes. It contains the numerator of the connected triples, which enter
the wave function at second order and the energy at fourth order, but no
triples denominator.

\subsubsection{Higher-Order Generalizations}

The same counting gives the leading order of every cross block. The adjoint of
a removal string of class $k$ fills $k+1$ holes and empties $k$ particles, so
it acts on ground-state components of excitation rank at least $k+1$; one
matrix element of $W$ changes the number of particles by at most two, so it
reaches an attachment configuration of class $l$ only from rank $k+l-1$ on;
the second term of Eq.~\eqref{eq:si:Xraw} gives the mirror condition. With the
orders at which these components first appear, the block between classes $k$
and $l$ starts no earlier than order
\begin{equation}
L(k,l)=\Big\lceil\frac{k+l+1}{2}\Big\rceil \;.
\label{eq:si:Lkl}
\end{equation}
This bound is our counting for the raw element; that the orthonormalization
does not lower it rests on the numerical checks. Under the order rule the
block between classes $k$ and $l$ is carried through order $n-k-l$, so the
class-1 block enters ADC(4), the blocks between classes 1 and 2 enter ADC(5),
and so on.

The two-sector and union-space ISR share all same-type blocks. This also
holds for the form of the two-sector ISR that builds the removal sector from
the removal strings themselves and takes the commutator with $H$ in reversed
order: for real orbitals each of its elements equals the corresponding
union-space element, for the configuration block
\begin{equation}
\langle\Psi_0|\{O^{-\dagger}_\mu,[O^-_{\mu'},H]\}|\Psi_0\rangle
=\braket{v(O^{-\dagger}_\mu)|\mathcal{H}|v(O^{-\dagger}_{\mu'})} \;,
\notag
\end{equation}
and likewise for the metric and the couplings. The difference of the two
folds therefore consists of the chains that contain at least one cross
block. A chain begins and ends with a coupling of at least first order,
and the cross block is at least of second order. The lowest such chain is
\begin{equation}
\pmb{\Delta}^{(4)}(\omega)=\mathbf{U}^{-(1)}(\omega+\mathbf{K}^-)^{-1}\mathbf{X}^{(2)}(\omega-\mathbf{K}^+)^{-1}\mathbf{U}^{+(1)\dagger}+\mathrm{h.c.} \;,
\label{eq:si:defect}
\end{equation}
the path from orbital $p$ through a 2h1p configuration, the cross block and a
2p1h configuration back to orbital $q$, with orders one, two and one. Every
other chain with a cross block is of fifth order or higher. The two-sector
ISR-ADC($n$) is therefore exact through third order for every $n$ and in
general not at fourth order. This is not a truncation effect: the complete
two-sector representation misses the same coefficient, because the deleted
block is part of the complete union-space matrix.

The cross block has a simple physical reading. The type of an intermediate
state is sharp only at zeroth order. In the correlated ground state each
intermediate state acquires a first-order admixture of the other sector,
Eqs.~\eqref{eq:si:admix1} and~\eqref{eq:si:admix2}, and the Hamiltonian couples
two states through these admixtures, at second order. The two-sector
construction keeps the effect of the admixtures on the same-type blocks but
drops their mutual coupling. In diagram language the missing content is a
3p3h fluctuation in flight: the energy difference between the two coupled
configurations is the excitation energy of $\ket{T_{\mu\nu}}$, and the
corresponding fourth-order time orderings carry the weight of a 3p3h times a
2p2h denominator. Part of this content is a ground-state triples amplitude, the
first term of Eq.~\eqref{eq:si:X2}. The rest divides matrix elements between
$(N\pm1)$-electron states by a triples energy and is not an amplitude of the
ground state, which is why no construction from the ground state alone
recovers it.\cite{tarantelli1992configuration}

%-----------------------------------------------------------------------------
\subsection{The Tarantelli--Cederbaum decoupling}
\label{sec:si:tc}

We now construct the transformation $\mathbf{V}$ of
Eq.~\eqref{eq:si:Vform} order by order.\cite{tarantelli1992configuration,Mertins_1996b} 
As we have discussed, block diagonalizations are not unique, as 
after decoupling any unitary transformation of the type Eq.~\eqref{eq:si:gauge} can be applied.

A canonical choice writes the configuration part of $\mathbf{V}$ as
the exponential of a generator $\mathbf{A}$ that connects only the two sectors,
\begin{equation}
\tilde{\mathbf{M}}'_{\mathrm{c}}=e^{-\mathbf{A}}\tilde{\mathbf{M}}_{\mathrm{c}}\,e^{\mathbf{A}}
=\tilde{\mathbf{M}}_{\mathrm{c}}+[\tilde{\mathbf{M}}_{\mathrm{c}},\mathbf{A}]
+\tfrac12\big[[\tilde{\mathbf{M}}_{\mathrm{c}},\mathbf{A}],\mathbf{A}\big]+\cdots \;,
\label{eq:si:bch}
\end{equation}
with $\mathbf{A}$ real antisymmetric, a power series in $\lambda$, and nonzero
only between a removal-type state $\mu$ and an attachment-type state $\nu$,
where we call its element $a_{\mu\nu}$.
This is the canonical van Vleck transformation of quasidegenerate perturbation
theory,\cite{VanVleck_1929,Shavitt_1980} which is also known as the
Schrieffer--Wolff transformation,\cite{Schrieffer_1966} here applied between
two blocks separated by a gap. It also has a closed form at finite
$\lambda$.\cite{Cederbaum_1989}

In the coefficient of order $m$ of Eq.~\eqref{eq:si:bch}, the new generator
$\mathbf{A}^{(m)}$ appears only through its commutator with the zeroth-order
diagonal, since every other term that contains it has an additional factor of
at least first order. On the diagonal stand minus $K^-_\mu$ and $K^+_\nu$, so
this commutator multiplies $a^{(m)}_{\mu\nu}$ by minus the sum of the two
energies. Let $\mathbf{R}^{(m)}$ be the cross part of the coefficient of order
$m$ computed from the lower orders of $\mathbf{A}$. The cross part vanishes if
and only if
\begin{equation}
a^{(m)}_{\mu\nu}=\frac{R^{(m)}_{\mu\nu}}{K^-_\mu+K^+_\nu} \;.
\label{eq:si:Arec}
\end{equation}
By Eq.~\eqref{eq:si:gap} every denominator is an excitation energy of rank
three or higher, so there are no small denominators. Since $\mathbf{X}$
vanishes at orders zero and one, the generator starts at second order with
\begin{equation}
a^{(2)}_{\mu\nu}=\frac{X^{(2)}_{\mu\nu}}{K^-_\mu+K^+_\nu} \;,
\label{eq:si:A2}
\end{equation}
the second-order triples coefficient plus the integral terms of
Eq.~\eqref{eq:si:X2} divided by the 3p3h energy. The recursion fixes every
order uniquely. Any other block diagonalization that reduces to the identity
at zeroth order differs from this one by a rotation within each sector: the
cross part of the connecting transformation obeys, order by order, an equation
whose zeroth-order operator is the difference of the attachment and removal
pole positions, which never vanishes, so it is zero. The other decouplings
are therefore exactly the gauge transformations of Eq.~\eqref{eq:si:gauge}.

\begin{table}[hbt!]
\caption{Blocks of the union-space ISR-ADC(4) with their lowest and highest orders. Same-type blocks occur once per sector, with the superscript $+$ or $-$ omitted.}
\label{tab:si:adc4}
\begin{tabular}{lcl}
\hline\hline
block & orders & content\\
\hline
$\pmb{\epsilon}+\pmb{\Sigma}(\infty)$ & 0--4 & orbitals\\
$\mathbf{U}_1$ & 1--3 & $\mathbf{U}^{(1)}_1+\mathbf{U}^{(2)}_1+\mathbf{U}^{(3)}_1$\\
$\mathbf{U}_2$ & 2 & $\mathbf{U}^{(2)}_2$\\
$\mathbf{K}_1+\mathbf{C}_{11}$ & 0--2 & $\mathbf{K}_1+\mathbf{C}^{(1)}_{11}+\mathbf{C}^{(2)}_{11}$\\
$\mathbf{C}_{12}$ & 1 & $\mathbf{C}^{(1)}_{12}$\\
$\mathbf{K}_2+\mathbf{C}_{22}$ & 0 & $\mathbf{K}_2$\\
$\mathbf{X}_{11}$ & 2 & $\mathbf{X}^{(2)}_{11}$\\
$\mathbf{X}_{12}$, $\mathbf{X}_{21}$ & -- & --\\
$\mathbf{X}_{22}$ & -- &  --\\
\hline\hline
\end{tabular}
\end{table}

We now work out ADC(4). Table~\ref{tab:si:adc4} lists the blocks of the
union-space ADC(4). The generator starts at second order and connects classes
1 and 1, 1 and 2, and 2 and 1. We go through the blocks and keep only what the
order rule keeps. The couplings to class 1 are carried through third order.
Their image under the generator is the first-order coupling times the
second-order generator, of third order, and it survives; all other images are
of fourth order or higher. The couplings to class 2 are carried through second
order, and their lowest image is of third order. The class-1 blocks are
carried through second order, while the first change of any same-type block is
of fourth order: the same-type part of Eq.~\eqref{eq:si:bch} arises only from
the cross block times the generator and from the double commutator, which
together give half the commutator of $\mathbf{X}^{(2)}$ with
$\mathbf{A}^{(2)}$. The remaining blocks change only at third order or later,
beyond what the rule carries. The cross blocks vanish by construction and the
orbital block is untouched.

The decoupled ADC(4) is therefore the two-sector ISR-ADC(4) with one change:
third-order rows are added to the class-1 couplings of both types,
\begin{align}
\delta\mathbf{U}^+_{\nu}&=\sum_{\mu}\mathbf{U}^{-(1)}_{\mu}\,
\frac{X^{(2)}_{\mu\nu}}{K^-_\mu+K^+_\nu} \;,
\notag\\
\delta\mathbf{U}^-_{\mu}&=-\sum_{\nu}\mathbf{U}^{+(1)}_{\nu}\,
\frac{X^{(2)}_{\mu\nu}}{K^-_\mu+K^+_\nu} \;,
\label{eq:si:deltaU}
\end{align}
where a coupling with a configuration subscript denotes its column of
first-order couplings to all orbitals and the sums run over class 1. These rows
are the partial fractions of the defect of Eq.~\eqref{eq:si:defect}. Splitting
the product of the two resolvents,
\begin{equation}
\frac{1}{(\omega+K^-_\mu)(\omega-K^+_\nu)}=\frac{1}{K^-_\mu+K^+_\nu}
\left[\frac{1}{\omega-K^+_\nu}-\frac{1}{\omega+K^-_\mu}\right] ,
\notag
\end{equation}
turns Eq.~\eqref{eq:si:defect} into single poles at class-1 configurations of
both types,
\begin{align}
\pmb{\Delta}^{(4)}(\omega)={}&\sum_\nu\frac{\delta\mathbf{U}^+_\nu\mathbf{U}^{+(1)\dagger}_\nu
+\mathrm{h.c.}}{\omega-K^+_\nu}
+\sum_\mu\frac{\delta\mathbf{U}^-_\mu\mathbf{U}^{-(1)\dagger}_\mu
+\mathrm{h.c.}}{\omega+K^-_\mu} \;,
\label{eq:si:defectpf}
\end{align}
which is exactly the first-order change of the two sector folds caused by the
new rows. The cross block couples a pole in one sector to a pole in the other,
and the partial fraction distributes this coupling over the two sectors with
the weight of a 3p3h energy. The three constructions keep this content in
three forms. The two-sector ISR deletes it. The decoupling places it in the
couplings, as the diagrammatic ansatz requires. The union-space ISR keeps the
block, and its eigenvalue problem generates the 3p3h denominators by coupling
the sectors. At fifth order the same counting shows that the decoupling still
changes only couplings; the first change of a configuration block within the
order rule appears in ADC(6).

%-----------------------------------------------------------------------------
\subsection{Proof for general $n$}
\label{sec:si:proof}

This section proves that the canonical decoupling of the union-space ISR,
truncated by the order rule, is a diagrammatic ADC($n$). We first fix the
objects and the one assumption the proof needs, then state the theorem, prove
four lemmas, derive the theorem from them, and close with its consequences and
limits.

\subsubsection{Definitions and assumption}

The proof compares matrices order by order, so we first say what it means for a
block to begin at a given order.

\begin{definition}
\label{def:si:start}
A block of a matrix whose elements are power series in $\lambda$ starts no
earlier than order $p$ if all its Taylor coefficients below order $p$ vanish.
\end{definition}

The target of the proof is the construction of Schirmer, Cederbaum and
Walter,\cite{Schirmer_1983} stated without reference to how the matrix elements
are found.

\begin{definition}
\label{def:si:twosector}
A two-sector representation of order $n$ is a static orbital matrix together
with an attachment fold and a removal fold of the form of the ansatz of
Sec.~\ref{sec:si:diag}. Its configuration spaces hold the classes of the order
rule, and its blocks are carried through the orders of the order rule.
\end{definition}

\begin{definition}
\label{def:si:diagadc}
A diagrammatic ADC($n$) is a two-sector representation of order $n$ with three
properties. Its attachment fold reproduces the exact attachment part of
Eq.~\eqref{eq:si:pmsplit} through order $n$ at every frequency. Its removal
fold does the same for the exact removal part. Its static matrix reproduces the
orbital energies plus the exact static self-energy through order $n$.
\end{definition}

The candidate is built from the complete union-space matrix of
Eq.~\eqref{eq:si:union} and the canonical generator of Sec.~\ref{sec:si:tc}.

\begin{definition}
\label{def:si:decadc}
The decoupled ADC($n$) is obtained in two steps. The configuration block of the
complete union-space matrix is transformed with the canonical generator of
Eqs.~\eqref{eq:si:bch} and~\eqref{eq:si:Arec}. The transformed matrix is then
truncated by the order rule, with the orbital block carried through order $n$.
\end{definition}

The canonical orders of Sec.~\ref{sec:si:isr} bound the same-type blocks. For
the cross blocks we have only the bound of Eq.~\eqref{eq:si:Lkl}, and the proof
needs it to be at least as strong as the class difference.

\begin{assumption}
\label{ass:si:A3}
For every removal-type class $k$ and every attachment-type class $l$, the cross
block between them starts no earlier than the class difference, that is,
\begin{equation}
L(k,l)\ \ge\ |k-l| \;.
\label{eq:si:A3}
\end{equation}
\end{assumption}

By Eq.~\eqref{eq:si:Lkl}, the assumption holds unless one class exceeds three
times the other plus two. It fails first for a removal-type class 1 and an
attachment-type class 6.

\subsubsection{Theorem}

\begin{theorem}
\label{thm:si:decoupling}
The decoupled ADC($n$) is a diagrammatic ADC($n$) for every $n$ from two to
nine. Under Assumption~\ref{ass:si:A3} it is a diagrammatic ADC($n$) for every
$n$.
\end{theorem}

The proof rests on two properties of the decoupled matrix. Its complete form
separates the sectors exactly (Lemma~\ref{lem:si:separation}), and its blocks
keep the canonical orders (Lemma~\ref{lem:si:orders}). Two general facts, the
uniqueness of the sector split (Lemma~\ref{lem:si:split}) and the effect of a
truncation on the fold (Lemma~\ref{lem:si:trunc}), connect these properties to
the definition of a diagrammatic ADC($n$).

\subsubsection{Lemmas}

The first lemma lets us match the two sectors separately. It restates the
argument of Sec.~\ref{sec:si:setting}.

\begin{lemma}[Split lemma]
\label{lem:si:split}
Let $f$ be a rational function of $\omega$ that vanishes at infinity and has
its poles in the union of $P^+$ and $P^-$. Then $f$ is the sum of exactly one
function with poles in $P^+$ and one function with poles in $P^-$, both
vanishing at infinity.
\end{lemma}

\begin{proof}
The partial-fraction decomposition of $f$ gives one such sum. The difference of
two such sums is a function whose poles lie both in $P^+$ and in $P^-$. The two
sets are disjoint, so this function has no poles. A rational function without
poles that vanishes at infinity is zero.
\end{proof}

The second lemma states which blocks a truncation may drop. It restates the
truncation argument of Sec.~\ref{sec:si:isr}.

\begin{lemma}[Truncation lemma]
\label{lem:si:trunc}
Assign to each class an integer weight, zero for the orbitals. Let a matrix of
the form of Eq.~\eqref{eq:si:union} have three properties. Every off-diagonal
block starts no earlier than the difference of the weights of its two classes.
Every coupling starts no earlier than the weight of its class. Every diagonal
block is the zeroth-order configuration energies plus terms of first or higher
order. Then the fold changes only beyond order $n$ if we drop the terms of each
block beyond order $n$ minus both weights, the terms of each coupling beyond
order $n$ minus its weight, and every class whose weight exceeds half of $n$.
\end{lemma}

\begin{proof}
By Eq.~\eqref{eq:si:trunc}, a chain of Eq.~\eqref{eq:si:chain} that passes
through a block term of order $p$ between classes of weights $d_I$ and $d_J$
has order at least $d_I$ plus $p$ plus $d_J$. A dropped block term has $p$
above $n$ minus both weights, and a dropped coupling term likewise. A chain that
enters an omitted class reaches it and leaves it at a total order of at least
twice its weight, which exceeds $n$. Every chain that contains a dropped term
is therefore of order $n+1$ or higher.
\end{proof}

The third lemma shows that the complete decoupled matrix already produces the
two exact sector parts.

\begin{lemma}[Sector separation]
\label{lem:si:separation}
The canonical decoupling of the complete union-space matrix has an attachment
fold and a removal fold. At every order $m$, the Taylor coefficient of the
attachment fold equals the exact attachment part of Eq.~\eqref{eq:si:pmsplit},
and that of the removal fold equals the exact removal part. Its orbital block
equals the orbital energies plus the exact static self-energy.
\end{lemma}

\begin{proof}
The complete union-space fold equals the exact dynamic self-energy by
Eq.~\eqref{eq:si:unionexact}. The decoupling does not touch the orbital
vectors, so the decoupled fold equals the union-space fold, and the orbital
block is that of Eq.~\eqref{eq:si:union}. The decoupled configuration block has
no cross part, so the decoupled fold is the sum of an attachment fold and a
removal fold, Eq.~\eqref{eq:si:Vfolds}. The generator starts at second order,
so the zeroth-order diagonals of the two folds are the attachment and the
removal pole positions. By the chain expansion, every Taylor coefficient of
each fold then has its poles in its own set. Lemma~\ref{lem:si:split}, applied
at each order, identifies the two folds with the exact attachment and removal
parts.
\end{proof}

Lemma~\ref{lem:si:separation} shows that the canonical decoupling of the exact
ISR produces, without a single diagram, the two parts that the Goldstone
expansion produces by sorting time orderings. The fourth lemma shows that the
decoupling keeps the orders that Lemma~\ref{lem:si:trunc} needs.

\begin{lemma}[Order preservation]
\label{lem:si:orders}
Under Assumption~\ref{ass:si:A3}, the decoupled complete matrix satisfies the
hypotheses of Lemma~\ref{lem:si:trunc} with the class index as weight.
\end{lemma}

\begin{proof}
In the union-space matrix, every same-type block starts no earlier than the
class difference by the canonical orders, and every cross block does so by
Eq.~\eqref{eq:si:Lkl} and Assumption~\ref{ass:si:A3}. We show by induction on
the order $m$ that the generator obeys the same bound. Suppose it does below
order $m$. The coefficient $R^{(m)}_{\mu\nu}$ of Eq.~\eqref{eq:si:Arec} is a
sum of products of blocks along index paths from $\mu$ to $\nu$ whose orders
add up to $m$. Each link starts no earlier than the class difference of its two
ends. By the triangle inequality, the coefficient vanishes when $m$ is smaller
than the class difference of $\mu$ and $\nu$. Equation~\eqref{eq:si:Arec}
divides it by a number that does not depend on $\lambda$, so the generator obeys
the bound at order $m$. The same triangle inequality carries the bound to the
exponential of the generator, to the transformed couplings and to the
transformed configuration block. The generator starts at second order, so the
diagonal of the transformed configuration block still starts with the
configuration energies, and every other diagonal term is of first or higher
order.
\end{proof}

\subsubsection{Proof of the theorem}

We first prove the theorem under Assumption~\ref{ass:si:A3} and then remove the
assumption for orders up to nine.

\begin{proof}[Proof of Theorem~\ref{thm:si:decoupling}]
By Lemmas~\ref{lem:si:orders} and~\ref{lem:si:trunc}, truncating the decoupled
complete matrix by the order rule changes its fold only beyond order $n$. The
truncation creates no cross blocks, so the truncated fold is still a sum of an
attachment fold and a removal fold, with Hermitian, frequency-independent
blocks over the classes of the rule. The decoupled ADC($n$) is therefore a
two-sector representation of order $n$. Its total fold agrees with the exact
dynamic self-energy through order $n$ by Lemma~\ref{lem:si:separation}.
Lemma~\ref{lem:si:split}, applied at each order through $n$, shows that each
sector fold agrees with its exact part. The static matrix is the orbital block,
which the decoupling does not touch and which equals the orbital energies plus
the exact static self-energy by Lemma~\ref{lem:si:separation}. This proves the
theorem under Assumption~\ref{ass:si:A3}.

Without the assumption, Eq.~\eqref{eq:si:Lkl} still bounds every cross block,
and Assumption~\ref{ass:si:A3} enters only Lemma~\ref{lem:si:orders}. Its first
failure is the block between a removal-type class 1 and an attachment-type
class 6. A 7p6h configuration can then be reached from the orbitals at fifth
instead of sixth order, and a chain through it and back contributes at tenth
order, before the order rule includes class 6. Every later violation acts at
still higher order. For $n$ up to nine, no chain affected by a violation
contributes through order $n$, and the argument above applies unchanged.
\end{proof}

\subsubsection{Consequences}

The theorem and its lemmas have three consequences for the constructions of
Table~\ref{tab:si:constructions}. The first concerns the order rule itself and
the union-space ISR.

\begin{corollary}[Order rule]
\label{cor:si:orderrule}
For every $n$ covered by Theorem~\ref{thm:si:decoupling}, the order rule of
Sec.~\ref{sec:si:diag} suffices for an ADC($n$). In particular, the union-space
ISR-ADC($n$) is exact through order $n$.
\end{corollary}

\begin{proof}
The argument of Lemma~\ref{lem:si:orders} for the union-space matrix itself
shows that it satisfies the hypotheses of Lemma~\ref{lem:si:trunc} with the
class index as weight. With this weight, Lemma~\ref{lem:si:trunc} is the order
rule. The complete union-space fold is exact by
Eq.~\eqref{eq:si:unionexact}.
\end{proof}

The second consequence concerns the principal poles, the quantities compared in
the main text.

\begin{corollary}[Principal poles]
\label{cor:si:poles}
The principal poles of the decoupled and of the union-space ADC($n$) have
Taylor series that agree with the exact ones through order $n$. Class-1
satellites are exact only through order $n-2$, the standard ADC result.
\end{corollary}

\begin{proof}
The Taylor coefficients of a principal pole are polynomials in the self-energy
coefficients through order $n$ and in their frequency derivatives at the
orbital energy, where the self-energy is regular. By
Theorem~\ref{thm:si:decoupling} and Corollary~\ref{cor:si:orderrule}, these
coefficients are exact.
\end{proof}

The theorem fixes the folds, not the matrices: the matching determines the
matrices only up to the gauge transformations of Eq.~\eqref{eq:si:gauge}. At
class 1 this is the only freedom, which gives the third consequence, an
element-wise statement at fourth order.

\begin{corollary}[Element-wise equality at ADC(4)]
\label{cor:si:adc4}
Let the class-1 configuration energies be nondegenerate and the first-order
couplings nonzero. Then any diagrammatic ADC(4) that uses the ISR's
second-order class-1 block has the same third-order class-1 couplings as the
decoupled ADC(4).
\end{corollary}

\begin{proof}
If two diagrammatic solutions agree through the lower orders, the residues of
their difference at each class-1 pole force the difference to be an
infinitesimal gauge transformation. With the class-1 block fixed at second
order, this transformation vanishes, so the third-order couplings coincide.
\end{proof}

The structural rule of Sec.~\ref{sec:si:diag} places the new fourth-order
content in exactly these couplings, since the time orderings that carry it have
a single frequency-dependent denominator.

\subsubsection{Remarks}

Three remarks delimit the result.

\begin{remark}
The theorem does not say that the decoupled ISR is the same matrix as one
obtained by sorting diagrams. Beyond fourth order we take the canonically
decoupled ISR as the definition of diagrammatic ADC($n$); any other
diagrammatic ADC($n$) differs from it by a gauge transformation.
\end{remark}

\begin{remark}
The bound of Eq.~\eqref{eq:si:Lkl} is our count for the raw cross element.
That the orthonormalization does not lower it rests on the numerical checks.
The unconditional part of Theorem~\ref{thm:si:decoupling}, for $n$ up to nine,
uses this bound as well.
\end{remark}

\begin{remark}
For $n$ of ten or more and systems with 7p6h configurations,
Lemma~\ref{lem:si:trunc} still applies if each class is weighted by the lowest
order at which any chain reaches it. Whether the extra blocks then contribute
we have not determined.
\end{remark}

%-----------------------------------------------------------------------------
\subsection{Non-Dyson ADC}
\label{sec:si:nondyson}

The construction so far concerns the Dyson self-energy.
Non-Dyson ADC\cite{Schirmer_1998,Trofimov_2005} needs only one sector. For
ionization, the intermediate states are built in the $(N-1)$-electron space
from the removal operators applied to the correlated ground state, with the
ordinary scalar product. With the matrix $\mathbf{M}$ of $H-E_0$ in these
states and the transition amplitudes $\mathbf{f}$ of the orbital removal
operators,
\begin{equation}
\mathbf{G}^-(\omega)=\mathbf{f}^\dagger(\omega+\mathbf{M})^{-1}\mathbf{f} \;,
\notag
\end{equation}
and the eigenvalues of $\mathbf{M}$ are the ionization energies. There is no
cross block, the same canonical orders and truncation lemma apply, and the
one-sector version of the proof shows exactness through order $n$. The
attachment sector enters only through the correlated ground state.

Table~\ref{tab:si:constructions} summarizes the four constructions of the main
text. They are built from the same intermediate states and differ only in what
they do with the cross content. All give finite Hermitian matrices, hence real
poles and positive residues. Beyond order $n$ they differ, and so do their
poles at the physical coupling. The ADC($n$) codes used in the main text were
produced as described in Sec.~\ref{sec:si:generator}.

\begin{table}[htb]
\caption{The four ADC constructions: treatment of the cross block
$\mathbf{X}$, order through which the self-energy (non-Dyson: the poles and
residues of one branch) is exact, and whether the complete construction is
exact.}
\label{tab:si:constructions}
\begin{tabular}{llcc}
\hline\hline
construction & cross block & exact & complete\\
\hline
diagrammatic (Dyson) & decoupled & $\lambda^n$ & yes\\
union-space ISR & kept & $\lambda^n$ & yes\\
two-sector ISR & deleted & $\lambda^3$ & no\\
non-Dyson ISR & none & $\lambda^n$ & yes\\
\hline\hline
\end{tabular}
\end{table}

% SI Part 2 -- the lambda oracle
\section{The $\lambda$-oracle}\label{sec:si:oracle}

The convergence data of Figs.~\ref{fig:NeConvergence} and \ref{fig:NeMP} of the main text, and the
corresponding data for Ar and Be, were obtained with one instrument that we call the $\lambda$-oracle.
It works in two separate steps.
The first step (Sec.~\ref{sec:si:oracle:exact}) computes exact quantities of a scaled Hamiltonian by
full configuration interaction (FCI): the IPs at the physical interaction strength, and the exact Taylor
coefficients of the ground-state energy, of the IPs and of the self-energy.
The second step (Sec.~\ref{sec:si:oracle:approx}) builds the approximations of
Fig.~\ref{fig:NeConvergence}, ADC($n$), GF($n$), IP-EOM-CC and IP-UCC, each from exact ingredients of the
first step together with the truncation rule that defines the method.
Both steps use the same orbitals, the same determinant spaces and the same Hamiltonian.
Every difference between a method and FCI in the figures is therefore the error of the method's
truncation and nothing else.

\subsection{Step A: the exact oracle}
\label{sec:si:oracle:exact}

All quantities refer to the Hamiltonian
\begin{equation}
H(\lambda) = F + \lambda\,(H - F) ,
\label{eq:si:Hlambda}
\end{equation}
in which $F$ is the canonical Hartree--Fock (HF) Fock operator and $H$ the physical Hamiltonian.
This is the M{\o}ller--Plesset partitioning, and we call $W$ the difference between $H$ and $F$.
The HF orbitals and orbital energies are those of the physical problem and are kept fixed for all
$\lambda$.
The zeroth-order Hamiltonian $F$ is diagonal in the basis of Slater determinants, with the sums of
orbital energies as eigenvalues, and for unit $\lambda$ the scaled Hamiltonian is the physical one.
A quantity ``of order $n$'' is the coefficient of $\lambda^n$ in its Taylor expansion, and a method is
``exact through order $n$'' if all its coefficients up to $\lambda^n$ are exact.
With this partitioning every diagram of order $n$ in the interaction carries exactly $\lambda^n$, while
all energy denominators keep their $\lambda$-independent values.
Frozen core orbitals (Ne 1s, the Ne-like core of Ar; none for Be) are doubly occupied in every
determinant and enter only through an effective one-electron operator.

The exact Green's function at a given $\lambda$ consists of an attachment and a removal branch,
\begin{align}
G_{pq}(\omega;\lambda) ={}& \braket{\Psi_0|a_p\big(\omega-H(\lambda)+E_0\big)^{-1}a^\dagger_q|\Psi_0}
\nonumber\\
&+ \braket{\Psi_0|a^\dagger_q\big(\omega+H(\lambda)-E_0\big)^{-1}a_p|\Psi_0} ,
\label{eq:si:or:G}
\end{align}
where $\ket{\Psi_0}$ and $E_0$ are the normalized FCI ground state and energy of $H(\lambda)$, and the two
resolvents act in the $(N+1)$- and $(N-1)$-electron determinant spaces.
The exact self-energy follows from the Dyson equation at each frequency,
\begin{equation}
\pmb{\Sigma}(\omega;\lambda) = \mathbf{G}_0^{-1}(\omega) - \mathbf{G}^{-1}(\omega;\lambda) ,
\label{eq:si:or:dyson}
\end{equation}
with $\mathbf{G}_0$ the HF propagator.
The FCI reference IP of each figure is the eigenvalue of $H$ in the $(N-1)$-electron space, minus the FCI
ground-state energy, of the cation state with the largest pole strength for removal of an electron from
the probed orbital.

FCI gives the self-energy at one value of $\lambda$, not as a series.
At a frequency off the real axis the self-energy is analytic in $\lambda$ near the origin, and its
Taylor coefficients are given exactly by Cauchy's integral formula,
\begin{equation}
\pmb{\Sigma}^{(m)}(\omega) = \frac{1}{2\pi i}\oint_{|\lambda|=r} \frac{\pmb{\Sigma}(\omega;\lambda)}{\lambda^{m+1}}\, d\lambda ,
\label{eq:si:cauchy}
\end{equation}
where the circle of radius $r$ must not enclose a singularity of the self-energy in the complex
$\lambda$ plane.
We evaluate the self-energy by FCI at $N_\lambda$ equally spaced points on this circle and discretize
the integral by the trapezoidal rule,\cite{Lyness_1967}
\begin{align}
\pmb{\Sigma}^{(m)}(\omega) \approx{}& \frac{1}{N_\lambda r^m}\sum_{j=0}^{N_\lambda-1}
e^{-2\pi i mj/N_\lambda} \nonumber\\
&\times\pmb{\Sigma}\big(\omega;re^{2\pi i j/N_\lambda}\big) .
\label{eq:si:contour}
\end{align}
The sum is a discrete Fourier transform, so one transform returns all orders below $N_\lambda$ at once,
without any fit.
Its only error is aliasing: the coefficient of order $m$ is contaminated by those of orders
$m+N_\lambda$, $m+2N_\lambda$, \dots, damped by powers of $r^{N_\lambda}$.
This error is negligible as long as the circle lies well inside the radius of convergence, which we
check by repeating the extraction on a second, smaller circle.
For complex $\lambda$ the Hamiltonian is complex symmetric rather than Hermitian, and all inner products
are evaluated without complex conjugation, which is the analytic continuation of the real-$\lambda$
expressions.
The contour is the definition of the exact coefficients against which everything else is checked.

The contour requires FCI in the $N$- and $(N\pm1)$-electron spaces at many complex $\lambda$, which is
affordable only for small systems.
For Ne, Ar and Be we obtain the same coefficients from Rayleigh--Schr\"odinger (RS) recursions, which
need only the action of $H$ on vectors.
Because $F$ is diagonal in the determinant basis, every zeroth-order resolvent is a division of each
determinant coefficient by an orbital-energy difference, as in the high-order MP calculations of Olsen
and coworkers.\cite{Olsen1996}
With the HF determinant $\ket{\Phi_0}$ as zeroth order and intermediate normalization, the order-$m$
corrections to the ground state and its energy are
\begin{align}
\ket{\Psi^{(m)}_0} ={}& R_0\Big[W\ket{\Psi^{(m-1)}_0}
- \sum_{l=1}^{m}E^{(l)}_0\ket{\Psi^{(m-l)}_0}\Big] ,
\label{eq:si:rs}\\
E^{(m)}_0 ={}& \braket{\Phi_0|W|\Psi^{(m-1)}_0} ,
\end{align}
where $R_0$ is the resolvent of $F$ with the HF determinant projected out.
The $E^{(m)}_0$ are the MP energy corrections.
Wherever the normalized ground state is needed, the series is divided by the power series of its norm.
The same recursion, started from the Koopmans determinant of the hole in orbital $p$, gives the RS
series of the connected cation state in the $(N-1)$-electron space, and the difference of the two
energy series is the exact Taylor series of the IP,
\begin{equation}
\mathrm{IP}_p(\lambda) = \sum_m \lambda^m \big[E^{N-1,(m)}_p - E^{(m)}_0\big] .
\label{eq:si:ipseries}
\end{equation}
The energy coefficients of Fig.~\ref{fig:NeMP} of the main text are the coefficients of these two
series.
Equation~\eqref{eq:si:ipseries} is also the criterion against which every method of Step B is checked:
a method that is exact through order $n$ must reproduce its coefficients through $\lambda^n$.

The Taylor coefficients of the self-energy at a fixed real frequency follow from the same recursions.
In the removal branch of Eq.~\eqref{eq:si:or:G} the operator to be inverted is a power series in
$\lambda$,
\begin{align}
A(\lambda) ={}& \omega + F - E^{(0)}_0 + \lambda\big(W - E^{(1)}_0\big) \nonumber\\
&- \sum_{j\ge2}\lambda^j E^{(j)}_0 \equiv \sum_j \lambda^j A_j ,
\end{align}
and so is the vector it acts on, the annihilator $a_p$ applied to the ground-state series, whose order-$m$
part we call $b_m$.
The order-$m$ part of the solution of the linear system is
\begin{equation}
x_m = A_0^{-1}\Big[b_m-\sum_{j=1}^{m}A_jx_{m-j}\Big] ,
\label{eq:si:resolventseries}
\end{equation}
which requires only the inverse of the diagonal $A_0$.
The attachment branch is obtained in the same way.
Projection onto the annihilated ground-state series gives the Taylor coefficients of $\mathbf{G}$, whose
zeroth order is $\mathbf{G}_0$.
The inverse of $\mathbf{G}$ is again a power series, and by Eq.~\eqref{eq:si:or:dyson} the self-energy
coefficients are
\begin{align}
\big(\mathbf{G}^{-1}\big)^{(k)} &= -\mathbf{G}_0^{-1}\sum_{j=1}^{k}\mathbf{G}^{(j)}\big(\mathbf{G}^{-1}\big)^{(k-j)} , \\
\pmb{\Sigma}^{(k)}(\omega) &= -\big(\mathbf{G}^{-1}\big)^{(k)}(\omega) .
\end{align}
They contain the static and the dynamic part together, and they have poles only at the zeroth-order
energies of the $(N\pm1)$-electron determinants, because only $A_0$ is inverted.
On small test systems these coefficients agree with those of the contour, Eq.~\eqref{eq:si:contour}.

Neither $H(\lambda)$ nor any operator used here couples different irreducible representations of the
abelian point group or different spin projections.
We therefore work in one symmetry block at a time: the totally symmetric $N$-electron space with equal
numbers of $\alpha$ and $\beta$ electrons, and the $(N\mp1)$-electron spaces of the symmetry of the
probed orbital with one $\beta$ electron removed or added.
For a degenerate level such as Ne 2p one component suffices.
The Hamiltonian is never stored as a matrix; it acts on vectors through a direct-CI
routine\cite{Knowles_1984,Olsen_1988} as implemented in PySCF.\cite{Sun_2018,Sun_2020}

\subsection{Step B: structured approximations}
\label{sec:si:oracle:approx}

Every approximation of Fig.~\ref{fig:NeConvergence} of the main text is built from exact ingredients
of Step A: the RS series of the ground state, the exact action of $H$ in the determinant spaces, and,
for GF($n$), the exact self-energy coefficients.
The only approximation is the truncation rule that defines each method.
All methods are evaluated at unit $\lambda$ in the same symmetry block as the FCI reference, and the IP is
the root with the largest weight on the one-hole configuration of the probed orbital.
For the methods designed to be exact through order $n$, the RS series of this root, computed from the
matrix series of the method, agrees with Eq.~\eqref{eq:si:ipseries} through $\lambda^n$.

\subsubsection{ADC($n$)}
\label{sec:si:oracle:nondyson}

The ADC($n$) curves of Fig.~\ref{fig:NeConvergence} of the main text are non-Dyson IP-ADC($n$), the
intermediate-state representation (ISR) of the $(N-1)$-electron space described in
Sec.~\ref{sec:si:cross}.
The oracle constructs this ISR exactly in the determinant space, as a power series in $\lambda$.
The precursor states are the ionizing operators of classes 0, 1, 2, \dots\ (1h, 2h1p, 3h2p, \dots)
applied to the RS ground-state series.
Inserting the series of $\ket{\Psi_0}$, $E_0$ and $H(\lambda)$ into the overlap and Hamiltonian matrices
of the precursors gives their Taylor coefficients exactly.
The class-wise Gram--Schmidt and L\"owdin steps of Sec.~\ref{sec:si:isr} are then carried out on these
power series: projections become series products, and the inverse square root of an overlap matrix
becomes a binomial series that terminates at every finite order.
The result is the Taylor series of the ISR matrix, whose eigenvalues, for the complete set of classes and
to all orders, are the exact IPs of $H(\lambda)$.
ADC($n$) applies the order rule of Sec.~\ref{sec:si:isr} to this series: it keeps the classes up to
$\lfloor n/2\rfloor$ and carries the block between classes $k$ and $l$ through order $n-k-l$.
Its IP is an eigenvalue of the summed matrix at unit $\lambda$; no frequency-dependent equation is solved.
ADC(1) is Koopmans' theorem.

The three Dyson variants of Sec.~\ref{sec:si:adc}, the union-space ISR, the two-sector ISR and
diagrammatic ADC($n$) obtained by the Tarantelli--Cederbaum decoupling of Sec.~\ref{sec:si:tc}, are built
with the same series machinery in the union of the $(N+1)$- and $(N-1)$-electron spaces.
Because the union space doubles the vector length, they are limited to moderate class spaces, and
Fig.~\ref{fig:NeConvergence} of the main text shows only non-Dyson ADC($n$).

\subsubsection{GF($n$)}
\label{sec:si:oracle:gfn}

GF($n$) [Eq.~\eqref{eq:GFn} of the main text], called MBGF($n$) by Hirata and
coworkers,\cite{Hirata_2024} is the exact self-energy with its Taylor series cut after order $n$,
\begin{equation}
\pmb{\Sigma}^{[n]}(\omega) = \sum_{k=2}^{n}\pmb{\Sigma}^{(k)}(\omega) ,
\end{equation}
evaluated at unit $\lambda$ with the exact coefficients of Step A.
The IP follows from the Dyson equation.
We follow the eigenvalue $e_p(\omega)$ of the orbital-energy matrix plus $\pmb{\Sigma}^{[n]}(\omega)$
whose eigenvector has the largest weight on the probed orbital $p$, and solve
\begin{equation}
\omega = e_p(\omega)
\label{eq:si:gfroot}
\end{equation}
for $\omega$, which is minus the IP.
The truncated self-energy has poles of increasing multiplicity at the zeroth-order configuration
energies of the block.
We search for the root between the nearest such pole below and the nearest pole above the orbital energy
of $p$, the interval in which the main line starts at zeroth order.
If there are several roots, we take the one with the largest residue.\cite{Hirata_2024}
A self-energy of resolvent form always yields real quasiparticle solutions, but the truncated series
loses this form from third order on.
When the difference between $e_p(\omega)$ and $\omega$ does not change sign anywhere in the interval,
Eq.~\eqref{eq:si:gfroot} has no real solution there; these are the missing GF($n$) points in
Fig.~\ref{fig:NeConvergence} of the main text.

\subsubsection{IP-EOM-CC}
\label{sec:si:oracle:cc}

IP-EOM-CC\cite{Stanton_1993a,Krylov_2008} is computed in the same symmetry blocks with the physical
Hamiltonian.
The cluster operator $T$ contains one amplitude for every determinant $\Phi_\mu$ of the $N$-electron block
with excitation rank from 1 to $r$, so that rank 2 is CCSD and rank 6 is CCSDTQPH.
Because $T$ acts directly on determinant-space vectors, its exponential is applied exactly, without a
commutator expansion.
The amplitudes solve
\begin{equation}
\braket{\Phi_\mu|e^{-T}He^{T}|\Phi_0} = 0
\end{equation}
for all these determinants, and the ground-state energy is the HF component of $He^{T}\ket{\Phi_0}$.
The cation space contains the determinants with up to $r$ holes (1h and 2h1p for CCSD), and the IP is
the eigenvalue of the similarity-transformed Hamiltonian in this space minus the CC ground-state energy.
In Fig.~\ref{fig:NeConvergence} of the main text the CC level of rank $r$ is drawn at order $2r-1$.

\subsubsection{IP-UCC}
\label{sec:si:oracle:ucc}

ADC($n$) and IP-UCC($n$) differ only in their intermediate states.\cite{Mertins_1996,Hodecker_2022}
In IP-UCC the ground state is a unitary transformation of the HF determinant,
\begin{equation}
\ket{\Psi_0} = e^{\sigma}\ket{\Phi_0} , \qquad \sigma = T-T^\dagger ,
\label{eq:si:ucc}
\end{equation}
and the intermediate states are the transformed cation determinants $e^{\sigma}\ket{\Phi_J}$, ordered in
the same classes as the precursors of ADC($n$).
These states are orthonormal at every $\lambda$, so no Gram--Schmidt or L\"owdin step is needed, and the
secular matrix is
\begin{equation}
M_{IJ} = \braket{\Phi_I|e^{-\sigma}\big(H(\lambda)-E_0\big)e^{\sigma}|\Phi_J} .
\label{eq:si:uccM}
\end{equation}

In the perturbative IP-UCC($n$), $\sigma$ is a power series in $\lambda$ whose order-$k$ part is built from
an excitation operator $T_k$ of every rank.
The $T_k$ are fixed order by order by the requirement that Eq.~\eqref{eq:si:ucc} reproduces the normalized
RS ground state.
At order $k$ this condition is linear in the new amplitudes,
\begin{equation}
T_k\ket{\Phi_0} = \ket{\Psi^{(k)}_0} - \big[e^{\sigma_1+\dots+\sigma_{k-1}}\ket{\Phi_0}\big]^{(k)} ,
\end{equation}
where the bracket denotes the order-$k$ part.
This is equivalent to solving the UCC amplitude equations order by order with no amplitude rank
dropped.
The matrix of Eq.~\eqref{eq:si:uccM} is then expanded in $\lambda$ and truncated by the same order rule as
ADC($n$).
At orders 2 and 3 this corresponds to the IP-UCC2 and IP-UCC3 schemes of the
literature.\cite{Hodecker_2022,Dempwolff_2022}

The non-perturbative IP-UCC instead solves a UCC ground state of finite rank at the given $\lambda$.
Here $T$ contains all excitations of rank 1 to $r$ (rank 2 is UCCSD), and the exponential is applied
exactly, without truncating the Baker--Campbell--Hausdorff expansion.
The amplitudes solve the generalized Brillouin conditions
\begin{equation}
\braket{\Psi_0|\big[H(\lambda),\tau_\nu-\tau_\nu^\dagger\big]|\Psi_0} = 0
\end{equation}
for every excitation operator $\tau_\nu$ of rank 1 to $r$, and the energy is the expectation value of
$H(\lambda)$.
With all ranks kept, this is the exact ground state.
The cation space is restricted to the classes up to $K$, and Eq.~\eqref{eq:si:uccM} is diagonalized in
this space with no truncation in order.
The resulting matrix is Hermitian and frequency independent, and the method has two parameters, the
rank $r$ of the ground state and the class limit $K$ of the cation space, and no perturbative input.

% SI Part 3 -- code generation
\section{Code generation}\label{sec:si:generator}

The ADC(2)--ADC(5) results of Fig.~3 of the main text were computed with code
that was generated, not written by hand. Two generators were built. The first
derives the intermediate-state representation (ISR) of Sec.~\ref{sec:si:adc}
symbolically and turns every matrix element into tensor contractions. All
calculations of Fig.~3 use it, including the diagrammatic ADC(4) and ADC(5),
which it obtains from the union-space ISR by the canonical decoupling. The
second is a diagram engine that builds the exact self-energy from Goldstone
time orderings without intermediate states. It located the content that the
two-sector ISR misses at fourth order, before the cross block $\mathbf{X}$ of
Eq.~(\ref{eq:ISRunion}) of the main text was identified as its source. The
production code does not use it.

\subsection{The ISR generator}\label{sec:si:isrgen}

\subsubsection{Matrix elements and Wick contraction}

The generator works in the union space, in which every operator adds one
electron. For the $(N+1)$ sector these are the attachment operators
$a^\dagger_a$, $a^\dagger_a a^\dagger_b a_j$, and so on; for the $(N-1)$
sector they are the adjoints of the removal operators, $a^\dagger_i$,
$a^\dagger_i a^\dagger_j a_b$, and so on. We denote any such operator string by
$Y_J$. Three kinds of elements are needed over the MP ground state
$\ket{\Psi_0}$: the configuration block, the coupling of an orbital to a
configuration, and the metric,
\begin{align}
M_{IJ} &= \bra{\Psi_0}\acomm{Y_I^\dagger}{\comm{H}{Y_J}}\ket{\Psi_0} \;,
\label{eq:si:gen-M}\\
U_{pJ} &= \bra{\Psi_0}\acomm{a_p}{\comm{H}{Y_J}}\ket{\Psi_0} \;,\\
S_{IJ} &= \bra{\Psi_0}\acomm{Y_I^\dagger}{Y_J}\ket{\Psi_0} \;.
\end{align}
The anticommutator lets one set of operators describe attachment and removal
together, and the commutator with $H$ removes the ground-state energy, so no
total energy appears. After orthonormalization these elements become the
blocks $\mathbf{U}^{\pm}$, $\mathbf{K}^{\pm}+\mathbf{C}^{\pm}$ and $\mathbf{X}$
of Eq.~(\ref{eq:ISRunion}) of the main text.

The Hamiltonian is split into the Fock operator $f$ and the fluctuation
potential $v$. For each piece $h$ the nested commutators are written out as
four operator products,
\begin{equation}
\acomm{Y_I^\dagger}{\comm{h}{Y_J}}
= Y_I^\dagger h Y_J - Y_I^\dagger Y_J h + h Y_J Y_I^\dagger - Y_J h Y_I^\dagger \;,
\end{equation}
and each product is passed to \texttt{pdaggerq},\cite{Rubin_2021,Liebenthal_2025}
which normal-orders it with respect to the Hartree--Fock (HF) determinant
$\ket{\Phi_0}$ and returns the fully contracted terms of Wick's theorem.

The correlated ground state enters through the MP wave operator $\Omega$ in
intermediate normalization. Expanding $\Omega$ into excitation operators
$t_k$ of rank $k$ and $\Omega^\dagger$ into de-excitation operators $l_j$,
every element of a given order becomes a finite sum of terms
\begin{equation}
\bra{\Phi_0} l_j\,\acomm{Y_I^\dagger}{\comm{h}{Y_J}}\,t_k\ket{\Phi_0} \;,
\label{eq:si:gen-split}
\end{equation}
divided order by order by the norm of $\ket{\Psi_0}$. We call each choice of
$l_j$, $h$ and $t_k$ a dressing split. Its operator content does not depend on
the perturbation orders of the two amplitudes, so one Wick contraction serves
every order at which the split contributes, and the amplitudes are inserted
as symbols. Doubles first appear at first order, singles at second order, and
an amplitude of rank three or higher at half its rank, rounded up; this fixes
which splits can contribute at each order. For the couplings at order $n$ the
highest surviving rank is one less than $n$, so ADC(5) needs amplitudes up to
rank four.

\subsubsection{Which blocks are needed}

The ISR blocks have definite leading orders.\cite{Schirmer_1991,Mertins_1996}
The coupling of the orbitals to class $\mu$ starts at order $\mu$, and the
block between classes $\mu$ and $\nu$ at the order given by their difference.
Every Taylor coefficient of the self-energy is a sum of chains of blocks from
the orbitals into the configuration space and back, and the order of a chain is
the sum of the orders of its links. This gives the ADC order rule: ADC($n$)
keeps the classes up to $n/2$ (rounded down) and carries the block between
classes $\mu_I$ and $\mu_J$ through order
\begin{equation}
n-\mu_I-\mu_J \;.
\end{equation}
Table~\ref{tab:si:gen-blocks} lists the result. The generator evaluates the
rule as a function of $n$; no block is listed by hand.

\begin{table}[htb]
\caption{Orders carried by the same-sector blocks of ADC($n$). Class 1 is the
$2h1p$ or $2p1h$ space, class 2 the $3h2p$ or $3p2h$ space; a dash means that
the block is absent.}
\label{tab:si:gen-blocks}
\begin{ruledtabular}
\begin{tabular}{lcccc}
block & ADC(2) & ADC(3) & ADC(4) & ADC(5)\\
\hline
orbitals--class 1 & 1 & 1--2 & 1--3 & 1--4\\
class 1--class 1 & 0 & 0--1 & 0--2 & 0--3\\
orbitals--class 2 & -- & -- & 2 & 2--3\\
class 1--class 2 & -- & -- & 1 & 1--2\\
class 2--class 2 & -- & -- & 0 & 0--1\\
\end{tabular}
\end{ruledtabular}
\end{table}

The same counting applies to the cross blocks. The class-1 cross block
$\mathbf{X}_{11}$ starts at second order,\cite{tarantelli1992configuration} so
the order rule first carries it in ADC(4). The cross blocks between class 1 of
one sector and class 2 of the other, $\mathbf{X}_{12}$ and $\mathbf{X}_{21}$,
also start at second order and first enter ADC(5).

\subsubsection{Orthonormalization and denominators}

The intermediate states are obtained from the precursor states
$Y_J\ket{\Psi_0}$ by orthogonalizing each class against the lower classes and
by symmetric orthonormalization within each class. Both steps act on power
series in $\lambda$ and need only products of the generated blocks, no further
Wick contraction. For the symmetric step, the series of the inverse
$\mathbf{Z}$ and of the inverse square root $\mathbf{Y}$ of the metric follow
order by order from the deviation $\Delta\mathbf{S}$ of the metric from the
unit matrix,
\begin{align}
\mathbf{Z}^{(m)} &= -\sum_{p=1}^{m}\Delta\mathbf{S}^{(p)}\mathbf{Z}^{(m-p)} \;,\\
\mathbf{Y}^{(m)} &= \frac12\Big[\mathbf{Z}^{(m)}
  -\sum_{p=1}^{m-1}\mathbf{Y}^{(p)}\mathbf{Y}^{(m-p)}\Big] \;,
\end{align}
both starting with the unit matrix at zeroth order. The orthonormalized
blocks are the order-$n$ parts of $\mathbf{U}\mathbf{Y}$ and
$\mathbf{Y}^\dagger\mathbf{M}\mathbf{Y}$.

Wick's theorem returns only numerators, products of Fock elements,
antisymmetrized integrals and amplitude symbols. The MP amplitude
denominators are reconstructed from the external indices of each amplitude,
and the zeroth-order energies $\mathbf{K}^{\pm}$ are orbital-energy sums with
the sign of the sector. The resolvent of Eq.~(\ref{eq:SigmaADC}) of the main
text never appears in the generated code: it is the eigenvalue problem solved
at runtime.

\subsubsection{Cross blocks and the diagrammatic construction}
\label{sec:si:gen-cross}

The cross element between a removal configuration $\mu$ and an attachment
configuration $\nu$ is Eq.~\eqref{eq:si:gen-M} with $Y_\mu$ as bra and $Y_\nu$
as ket, so the same machinery produces it. It has no terms at zeroth and
first order, and the cross metric vanishes at every order. At second order
only one dressing split survives, in which the first-order de-excitation
doubles amplitude is contracted with the two-electron integrals.

Three Dyson constructions follow, identical through ADC(3). The two-sector ISR
deletes $\mathbf{X}$. The union-space ISR keeps it as a block of the
supermatrix. The diagrammatic construction removes it by the canonical
decoupling of Tarantelli and Cederbaum\cite{tarantelli1992configuration} and
keeps its images through order $n$; Sec.~\ref{sec:si:adc} shows that this is
diagrammatic ADC($n$). The decoupling is an orthogonal transformation
$e^{\mathbf{A}}$ of the configuration space, with $\mathbf{A}$ antisymmetric and
nonzero only between the two sectors. Requiring the cross part of the
transformed configuration block to vanish order by order gives
\begin{equation}
A^{(m)}_{\mu\nu}=\frac{R^{(m)}_{\mu\nu}}{K^-_\mu+K^+_\nu} \;,
\label{eq:si:gen-Arec}
\end{equation}
where $R^{(m)}$ is the order-$m$ cross part computed with the lower orders of
$\mathbf{A}$; at second order it is $\mathbf{X}^{(2)}$. The denominator is the
energy of a $3h3p$ or higher excitation, so there are no small denominators.

Within the orders kept by the order rule, the decoupling changes only the
couplings through ADC(5). For ADC(4) the single surviving image is a
third-order row added to each class-1 coupling,
\begin{align}
\delta U^{+}_{p\nu} &= \sum_{\mu}\frac{U^{-(1)}_{p\mu}X^{(2)}_{\mu\nu}}{K^-_\mu+K^+_\nu} \;,
\label{eq:si:gen-du-plus}\\
\delta U^{-}_{p\mu} &= -\sum_{\nu}\frac{U^{+(1)}_{p\nu}X^{(2)}_{\mu\nu}}{K^-_\mu+K^+_\nu} \;,
\label{eq:si:gen-du-minus}
\end{align}
with $\mu$ running over $2h1p$ and $\nu$ over $2p1h$ configurations. For
ADC(5) the class-1 couplings receive in addition the fourth-order images of
$\mathbf{A}^{(2)}$ and $\mathbf{A}^{(3)}$, and the class-2 couplings the
third-order image of $\mathbf{A}^{(2)}$ built from $\mathbf{X}_{12}$ and
$\mathbf{X}_{21}$. In the ADC(5) calculations the division by the
denominator of Eq.~\eqref{eq:si:gen-Arec} was evaluated by a Laplace
quadrature,\cite{Almlof_1991,Takatsuka_2008,Kaltak_2014} a numerical
approximation; at ADC(4) the decoupling is exact.

\subsubsection{Static part and non-Dyson ADC}

The orbital block of the Dyson supermatrix contains $\pmb{\Sigma}(\infty)$,
obtained as in Ref.~\citenum{Keizer_2026} by contracting the antisymmetrized
integrals with the difference between the unrelaxed MP($n-1$) and the HF
one-particle density matrices. For each ionization the root with the largest
pole strength is taken.

Non-Dyson ADC uses the same machinery with the ordinary metric of
$(N-1)$-electron states instead of the anticommutator, only the $(N-1)$ sector,
and occupied annihilators as class 0. The generated elements are then the ISR
matrix and the transition amplitudes of Sec.~\ref{sec:si:adc}. Non-Dyson
ADC(5) needs the ground state through fourth order.

\subsection{The diagram engine}\label{sec:si:diagrams}

The diagram engine builds the exact self-energy order by order in $\lambda$
without intermediate states. It was built to find which part of the exact
fourth-order self-energy the two-sector ISR misses, and it answered this
before the cross block was identified as the missing piece. The cross-block
decoupling of Sec.~\ref{sec:si:gen-cross} then replaced the diagram-based
corrections; the blocks the engine emits are kept as optional dressings and
independent checks.

\subsubsection{Skeletons, time orderings and coefficients}

A skeleton is a self-energy diagram drawn as a graph: $n$ two-body vertices,
particle lines (upward in time), hole lines (downward), and two marked vertices
that carry the external orbitals. The engine enumerates all connected,
one-particle irreducible skeletons with at least one valid time ordering,
excluding lines that start and end on the same vertex, which cancel against
the Fock operator for a HF reference, and removes duplicates by a canonical
form of the graph. Generating graphs directly replaced an earlier route that
recognized diagrams in the \texttt{pdaggerq} expansion of a product of $n$
interactions, which produces each diagram many times and exceeded the memory
of a single machine at fifth order.

The time orderings of a skeleton are the arrangements of its vertices that
respect the line directions. Each is evaluated with the Goldstone rules: a
horizontal cut between consecutive vertices contributes the denominator
\begin{equation}
D_{\mathrm{cut}} = \pm\omega + \sum_{i\,\in\,\mathrm{cut}}\epsilon_i
                 - \sum_{a\,\in\,\mathrm{cut}}\epsilon_a \;,
\end{equation}
where the frequency appears only if the cut lies between the two external
vertices (an $\omega$-cut). Every other cut (a free cut) carries a
ground-state excitation energy. The value of an ordering is the product of its
antisymmetrized integrals divided by its cut denominators, times the
coefficient of the skeleton,
\begin{equation}
\mathrm{coeff} = \sigma\,(-1)^{h+\ell}\,
\Big(\frac12\Big)^{n_{\mathrm{ep}}}\,\frac{1}{|\mathrm{Aut}|} \;.
\end{equation}
Here $h$ is the number of hole lines, $\ell$ the number of closed loops,
$n_{\mathrm{ep}}$ the number of pairs of equivalent lines, $|\mathrm{Aut}|$ the
number of permutations of internal vertices that map the diagram onto itself,
and $\sigma$ a sign that distinguishes a virtual from an occupied external
orbital. The same enumeration reproduces the published numbers of energy
diagrams of high-order perturbation theory for nuclei through sixth
order.\cite{Li_2026}

\subsubsection{Cuts, quartets and the fourth-order defect}

An $\omega$-cut crossed by three lines gives a pole at a class-1 energy, one
crossed by five lines a pole at a class-2 energy. At fourth order some
orderings have an $\omega$-cut crossed by seven lines, a configuration with no
counterpart in the ADC(4) space. These orderings come in groups of four that
share the seven-line cut, and within each group the residues at the seven-line
pole cancel.\cite{Schirmer_1983} The engine sums each group in closed form, so
the spurious poles never appear numerically.

Because every cut denominator is linear in $\omega$, the residue of each
ordering at a class-1 pole has a closed form. Sorting the residues by which
cut vanishes at the pole, the engine paired every class of fourth-order
orderings with the ISR term that reproduces it, with one exception: the
orderings whose middle cut is free and crossed by six lines, a $3h3p$
configuration in flight between the two halves of the diagram. They have a
single $\omega$-cut and are therefore single poles at class-1 energies.
Removing the external vertex that carries the first-order coupling leaves a
static row per class-1 configuration with the $3h3p$ energy as denominator.
These rows are identical to the decoupled cross-block rows of
Eqs.~\eqref{eq:si:gen-du-plus} and \eqref{eq:si:gen-du-minus}: the engine and
the canonical decoupling reach the same object from two sides. At fifth order
the engine was used to fit further corrections, which the canonical
decoupling has replaced.

%%%%%%%%%%%%%%%%%%%%%%%%%%%%%%%%
% \section*{References}  % revtex adds its own References heading (and TOC entry)
%%%%%%%%%%%%%%%%%%%%%%%%%%%%%

%%%%%%%%%%%%%%%%%%%%%%%%
\bibliography{\maintextdir biblio}
%%%%%%%%%%%%%%%%%%%%%%%%

\end{document}

